% Français 12357, batch 2 — Gallica views 21–40.
% Pass: Opus 5.5 (claude-opus-5-5), 2026-09-22 — first pass, unchecked against
% the views by a human.
% Read from the local mirror (archives/fr-12357), in four full-width bands per
% view at 2400 px and tight crops for doubtful words, numbers and marginalia;
% each view was drafted to disk as it was read and the file assembled from the
% drafts. No edition of Mersenne and no copy of the Harmonie universelle was
% opened.
%
% What the views are. Greyscale scans, portrait, about 3920 × 5745, each showing
% ONE page of the bound volume, with a strip of the facing page at the edge.
% Rectos (even views 22–40, except 32 and 34) and versos (odd views) alternate.
% The text runs into the gutter on the versos, and many line-ends there are
% lost: \add{} completes a word only where it is certain, \ill{} stands
% otherwise.
%
% What the volume is. A COPY, in one hand, of Mersenne's marginal and flyleaf
% remarks in his own Harmonie universelle — the copyist says « ainsy qu'il les
% avoit escrites ». The copy keeps Mersenne's first person (« j'ay donné ma
% methode », « nostre version », « J'ay veu graver a Mons ») ; the leaves are
% not in Mersenne's hand. Each remark opens with a key to the printed book,
% written on a line of its own and usually separated from the preceding remark
% by a short rule in the left margin: « p. 113. au commencemt. », « p. 138
% ligne 21 apres ces mots (… ) », « p. 26. prop. 6. lig. 14. », « lig. 6.
% devant la fin », « ibid. prop. XVIII », « p. 42 au 3me advertissemt. lig. 5. ».
% They are transcribed exactly as written, on their own line, and are the keys
% a reader needs; none was checked against the book. The page numbers restart
% at each treatise of the book, and the copyist marks the change by a running
% title in the left margin, under a rule; those titles are given as \marginal.
%
% The hand. One regular, rapid French cursive of the seventeenth century, the
% same from view 21 to view 40, with long descenders and a looped d; the
% long s is set s. Abbreviations kept as written: pō (pour), coe (comme), ē
% (est), fe (faire), pr, pre, pers (premier, premiere, premiers), ptie (partie),
% -emt (-ement), tousrs (tousiours), mr, sr, nre (nostre), aure (autre), no
% (nous), propoon, lesd., susd., dud., cæt. (in « et cæt. »), soe (somme),
% R. (racine), and the pound sign, two linked t, set ^{tt}. Figures and musical
% examples are drawn by the same pen; they are described in \note{}, not drawn.
% Fractions are set \frac; « b fa ♮ mi » and the square b are set \natural.
%
% What the twenty views hold, by the running titles the copy gives.
%   v. 21–24 (top)  The end of a run of remarks on combinations, paginated
%                   p. 113–154 of their treatise: the songs of eight notes
%                   and the removing of bad intervals (40320 songs), the cards
%                   of piquet (12 cards out of 36: 599555620984320000 in order,
%                   1251677700 without), the 100 fruits chosen from 6 baskets,
%                   triangular powers, the general combination, prop. 12, 17,
%                   20, with references to the Latin Harmonie (« pag. 136 de
%                   l'harmonie en latin ») and to « mr. Fren. » (v. 23). A
%                   remark « p. 139 » corrects a figure of the printed text
%                   (3514140 « et non 3513840, coe il y a dans le texte »).
%   v. 24–28        Running title « … 1er des …sonances » (v. 24, partly in the
%                   gutter), then « Du livre 3me des consonnances » (v. 25):
%                   remarks p. 26–315 — the coincidence of vibrations (Galileo's
%                   first dialogue, Kircher), the division of the octave and of
%                   the fourth into superparticular ratios (a table of the 27
%                   divisions of the fourth, v. 25), the means, the tetrachords
%                   and the great system, Le Maire's seven syllables and his
%                   characters (v. 26), Jean Le Fevre's scales after Gafurius
%                   with tables of lengths, squares and weights (v. 26), the
%                   modes after Doni (v. 27, with staves), note values.
%   v. 28–31        « Des Instruments …chorde » (v. 28), « …ivre second des
%                   Instruments » (v. 29), « …ivre 3me des …truments a chorde »
%                   (v. 30): strings of gut and metal, a table of 24 equal
%                   quarter-tones (v. 29), the archlute, Le Maire's almerie, a
%                   lute tablature (v. 29), fret placing, woods, Faber's
%                   weights, Galileo on gold and brass strings, Faber's
%                   « musique de bois » with a table (v. 30), the monochord and
%                   Faber's division of it, Cæsarianus and Gafurius, Salinas,
%                   the harmonic cane of Cambaye, Faber on flutes (v. 31).
%   v. 32           A fold-out leaf of figures, photographed open: the two
%                   monochord scales of v. 31 and, headed « Livre 5me des
%                   instruments a vent propoon 6me p. 232 », a cornet, a fife,
%                   a recorder and a flageolet.
%   v. 33           Flutes (continued), wind instruments; running title
%                   « … vijme livre instruments de …ussion » (the numeral
%                   doubtful): the proportions of a bell.
%   v. 34           A small loose sheet with the drawing of the bell and its
%                   construction, lettered; no text.
%   v. 35–37        The construction and measures of the bell (two tables,
%                   v. 35), « l'aveugle d'Estampes » and « un aveugle
%                   d'Utrecht », the weight of bell-metal, glasses (Kircher),
%                   the bow, the archimole.
%   v. 37–39        « De l'utilité de l'harmonie » (v. 37): weighing air
%                   (Galileo, « Mr. des Cartes » 1 to 145), burning mirrors and
%                   telescopes with two figures (v. 38), rational ellipses for
%                   Descartes's Dioptrique with an 18-digit number (v. 39),
%                   Galileo's projectiles, ranges of arquebus and canon.
%   v. 40           « … nouvelles …servations …siques et mathematiques »:
%                   weighing air, falling bodies, refraction, Cardan, jets of
%                   water; stops mid-sentence at the foot of the page.
%
% A date on the leaf. View 39: « le dernier jr de may 1636. sur la Foire a St.
% Clou. », heading an account of arquebus ranges; the first digit is small and
% the year is set \uncertain.
%
% Blank views. None in the batch. View 32 (fold-out of figures) and view 34
% (loose sheet with the bell) carry no text but figures and a heading, and are
% kept with a descriptive \note.
%
% Numbers on the leaves. The only series is an ink figure at the top right of
% each recto, the copy's own numbering, one per leaf: 11 (v. 22), 12 (24),
% 13 (26), 14 (28), 15 (30), 16 (33), 17 (36), 18 (38), 19 (40). It is
% recorded in the \note{} of each view and never in \folio{}. The versos carry
% no number. No pale pencil foliation was seen on any recto, so no \folio{} is
% written. Other figures: an isolated ink « 2 » at the foot of v. 24 (perhaps a
% gathering mark); an ink « 81 » at the head of v. 39 that is, read in a
% mirror, the « 18 » of the recto seen through the leaf. All other numbers
% (« p. 113 », « lig. 14 », « pag. 136 de l'harmonie en latin », « p. 831 de
% magnetismo musica », « Faber pag. 123 ») are references, to the Harmonie
% universelle or to other books, and are content.
\documentclass[11pt,a4paper]{article}
% The preamble shared by every transcription of the mathematicians' manuscripts in Gallica.
%
% It rests on one idea: **uncertainty compounds, it does not resolve.** A
% transcription that smooths over a doubtful word has destroyed the only thing
% it was made for — a reader must be able to tell, at every point, what was on
% the page and what was supplied. The apparatus macros below are therefore the
% critical apparatus, not decoration.
%
% The same commands are recognised by `scripts/render.mjs`, which produces the
% site's reading view, and by `scripts/tei.mjs`. A macro added here must be
% added there too, or rendering fails loudly — which is the wanted behaviour.
%
% Comments here are English, like the rest of the repository. The documents
% this preamble sets are French, and that is deliberate: see the skills.

\ProvidesPackage{germain}[2026/09/15 Transcriptions of mathematicians' manuscripts in Gallica]

\usepackage{amsmath,amssymb,amsthm}
\usepackage{mathrsfs}
% Kept from the parent preamble so the renderer's diagram support stays
% available; these manuscripts draw no commutative diagrams, and nothing here uses it.
\usepackage{tikz-cd}
\usepackage[english,french]{babel}
\usepackage{fontspec}

% --- Greek ------------------------------------------------------------------
% The text face stays fontspec's default, Latin Modern, which has no Greek:
% XeTeX drops every glyph it lacks with a line in the log and nothing on the
% page, and the Greek words the manuscripts quote vanished from the PDFs. So
% the Greek and Greek Extended blocks get a XeTeX character class of their own,
% and entering or leaving a run of it switches the family to CMU Serif — the
% same Computer Modern design, with polytonic Greek — and back. Only the
% family changes: a Greek word in \textit or \textbf keeps its shape and series.
% The transitions to and from every other class carry the tokens that class
% has towards an ordinary letter, so French spacing before « ; » or inside
% « » still holds after a Greek word. No grouping, so no brace or \endgroup
% can come out unbalanced; maths is left alone, since XeTeX inserts nothing
% there. Kept identical in `exercices.sty`.
\makeatletter
\newfontfamily\germainGreekfont{cmun}[
  Extension=.otf, NFSSFamily=germaingreek,
  UprightFont=*rm, ItalicFont=*ti, SlantedFont=*sl,
  BoldFont=*bx, BoldItalicFont=*bi, BoldSlantedFont=*bl]
\newXeTeXintercharclass\germain@greek
\def\germain@greekrange#1#2{%
  \@tempcnta=#1\relax
  \loop
    \XeTeXcharclass\@tempcnta=\germain@greek
    \ifnum\@tempcnta<#2\advance\@tempcnta\@ne\repeat}
\germain@greekrange{"0370}{"03FF}
\germain@greekrange{"1F00}{"1FFF}
\def\germain@greekprev{\rmdefault}
\def\germain@greekin{%
  \edef\germain@greekprev{\f@family}\fontfamily{germaingreek}\selectfont}
\def\germain@greekout{\fontfamily{\germain@greekprev}\selectfont}
\def\germain@greekjoin{%
  \XeTeXinterchartokenstate=\@ne
  \@tempcnta=\z@
  \loop
    \ifnum\@tempcnta=\germain@greek\else
      \XeTeXinterchartoks\germain@greek\@tempcnta=\expandafter{%
        \expandafter\germain@greekout\the\XeTeXinterchartoks\z@\@tempcnta}%
      \XeTeXinterchartoks\@tempcnta\germain@greek=\expandafter{%
        \the\XeTeXinterchartoks\@tempcnta\z@\germain@greekin}%
    \fi
    \ifnum\@tempcnta<4095\advance\@tempcnta\@ne\repeat}
% babel-french sets its own classes, and saves and restores the interchar
% state, each time a language is selected: join again after it does.
\addto\extrasfrench{\germain@greekjoin}
\addto\extrasenglish{\germain@greekjoin}
\AtBeginDocument{\germain@greekjoin}
\makeatother

\usepackage{geometry}
\geometry{a4paper,margin=25mm,marginparwidth=28mm}
\usepackage{marginnote}
\usepackage{xcolor}
\usepackage[normalem]{ulem}
\setlength{\emergencystretch}{3em}
\usepackage{hyperref}
\hypersetup{colorlinks=true,linkcolor=black,urlcolor=black,pdfborder={0 0 0}}

% --- Batch metadata ---------------------------------------------------------
% Taken as they stand by the rendering script, which shows them at the head of
% the reading view. They fix what the file claims to cover. The macro names are
% the parent project's, so that the scripts stay shared; \folder here means the
% volume's slug — `fr-9115` — and \pages counts Gallica views.

\newcommand{\folder}[1]{\gdef\grothFolder{#1}}
\newcommand{\batch}[1]{\gdef\grothBatch{#1}}
\newcommand{\pages}[2]{\gdef\grothFirst{#1}\gdef\grothLast{#2}}
\newcommand{\foldertitle}[1]{\gdef\grothFoldertitle{#1}}

% The holder's own shelfmark, printed as they write it: « Français 9115 ».
\newcommand{\shelfmark}[1]{\gdef\grothShelfmark{#1}}

% The dating, copied verbatim from the holder's catalogue — never inferred
% here. « XIXe siècle » is the BnF's; a pass that dates a leaf from its content
% says so in a \note{}, as its own inference.
\newcommand{\dating}[1]{\gdef\grothDating{#1}}

% The status notice, printed in the title block and repeated as a diagonal
% watermark on every page of the PDF. The manuscripts are in the public
% domain; what this line declares is not a right but a quality — an
% unverified first pass by a machine. The skills write the line; a file
% without it gets no watermark, so leaving it out is visible in the output.
\newcommand{\watermark}[1]{\gdef\grothWatermark{#1}}

\gdef\grothFolder{?}\gdef\grothBatch{}\gdef\grothFirst{?}\gdef\grothLast{?}%
\gdef\grothFoldertitle{}\gdef\grothShelfmark{}\gdef\grothDating{}\gdef\grothWatermark{}

% --- The résumé -------------------------------------------------------------
\newenvironment{resume}
  {\small\setlength{\parskip}{0.45em}}
  {\par\medskip\hrule\medskip}

% --- Keywords ---------------------------------------------------------------
% In English, closing the modernised reading's résumé: the modern vocabulary
% under which to search for what the leaves construct. The manifest extracts
% them to tag the volume — this line is the single source of the tags.
\newcommand{\keywords}[1]{\par\smallskip\noindent
  {\footnotesize\textsc{Keywords} --- \textit{#1}}\par}

% --- The critical apparatus -------------------------------------------------

% The Gallica view number, in the margin. This is the anchor that lets a
% disagreement over a formula be settled by looking at *one* image rather than
% a volume of 750. The site uses it to turn the facsimile as the transcription
% is scrolled. A view is one scanned image: usually one side of a leaf.
\newcommand{\page}[1]{%
  \marginnote{\footnotesize\color{gray!70}\textsf{v.\,#1}}%
  \label{page:#1}\ignorespaces}

% The BnF's pencilled foliation, where the leaf carries it and a pass can read
% it: « 198r ». This is what the literature cites — « ff. 198r–208v » — and
% the one line that turns a citation into a view. Recorded, never inferred:
% a folio a pass cannot read is not written.
\newcommand{\folio}[1]{%
  \marginnote{\footnotesize\color{gray!50}\textsf{f.\,#1}}\ignorespaces}

% The modernised reading's coarser counterpart to \page{}: which views a
% section is drawn from.
\newcommand{\pagerange}[2]{%
  \marginnote{\footnotesize\color{gray!70}\textsf{v.\,#1--#2}}%
  \ignorespaces}

% Illegible: never guessed.
\newcommand{\ill}{\textcolor{red!60!black}{[\,\ldots\,]}}

% A doubtful reading: the word is given, and flagged as uncertain.
\newcommand{\uncertain}[1]{\uline{#1}}

% An editorial addition — a missing letter, an implied word.
\newcommand{\add}[1]{\textcolor{blue!50!black}{[#1]}}

% Struck out by the writer. What was crossed out is part of the document.
\newcommand{\struck}[1]{\sout{#1}}

% The transcriber's note, on the state of the page or the mathematics.
\newcommand{\note}[1]{\par\smallskip{\small\color{gray!60!black}\textit{[#1]}}\par\smallskip}

% A marginal note of hers, distinct from the body of the text.
\newcommand{\marginal}[1]{\par\smallskip\noindent\hspace{1em}%
  {\small\color{gray!60!black}#1}\par\smallskip}

% --- Title ------------------------------------------------------------------
% The title block is French, like the documents themselves.

\AddToHook{shipout/background}{%
  \ifx\grothWatermark\empty\else
    \begin{tikzpicture}[remember picture, overlay]
      \node[rotate=52, scale=3.4, align=center, text=gray!18,
            font=\sffamily\bfseries]
        at (current page.center) {\grothWatermark};
    \end{tikzpicture}%
  \fi}

\AtBeginDocument{%
  \begin{center}
    {\large\bfseries Manuscrits de mathématiciens (Gallica) ---
      \ifx\grothShelfmark\empty\grothFolder\else\grothShelfmark\fi}\\[2pt]
    {\itshape\grothFoldertitle}\\[4pt]
    {\small vues \grothFirst--\grothLast
      \ifx\grothBatch\empty\else\ (lot \grothBatch)\fi}\\[2pt]
    \ifx\grothDating\empty\else
      {\small\color{gray!60!black}Datation du catalogue : \grothDating}\\[2pt]
    \fi
    \ifx\grothWatermark\empty\else
      {\small\bfseries\color{red!45!gray}\grothWatermark}\\[2pt]
    \fi
    {\footnotesize\color{gray!60!black}Transcription établie d'après les images IIIF de Gallica.
     Source gallica.bnf.fr / Bibliothèque nationale de France.\\
     Les lectures incertaines sont soulignées, les illisibles notées [\,\ldots\,],
     les ajouts entre crochets ; \textsf{v.} numérote les vues de Gallica,
     \textsf{f.} la foliotation au crayon de la BnF.}
  \end{center}
  \vspace{1em}}

\endinput


\folder{fr-12357}
\shelfmark{Français 12357}
\batch{2}
\pages{21}{40}
\dating{1601-1700}
\watermark{Lecture automatique\\non vérifiée}
\foldertitle{« Remarques tirées du livre de l'Harmonie universelle du P. MERSENNE, ainsy qu'il les avoit escrites, de sa main, à la marge et aux feuillets blancs devant et derrière dudit livre ».}

\begin{document}

\page{21}
\note{Page de gauche (verso) ; aucun chiffre en tête. Le texte court jusque dans le pli de la reliure : les dernières lettres de la plupart des lignes y sont cachées ; elles sont complétées entre crochets par \add{} seulement quand le mot est certain.}

p. 113. au commencem\textsuperscript{t}.

pō scauoir comme il faut separer les mauuais chants, ou telz autres interualles qu'on n\ill{} de ceste combinaison, par exemple pō oster vt bi, ou bien vt la, des 4320 chants de l'octa\add{ve} prenez vt, bi pō vne note, dont la combinaison auec les 6 autres notes sera comme la combinaison de 7 notes, c'est a dire 5040, et parce que bi, vt, donne encore le mesme interv\add{alle} prenez le double scauoir 10080: de tout autant pō vt, la et la, vt. mais parce que parmy ces 20160 chants, il s'en trouue quelques vns qui ont ensemble vt bi, et la, \ill{} scauoir vne septiesme maieure, et vne sextiesme; il faut prendre bi, vt la pō vne n\add{ote} qui auec les 5 notes fait 6 notes, on aura la combinaison de 6 choses, a scauoir 720: et pource qu'on peut aussy prendre la, vt, bi. Il faut doubler ce nombre, pō auoir 1440, le\add{quel} osté de 20160, reste 18720 chants, qui ont vne septe maieure, ou vne sextiesme qu'il faut oster de 40320. Il faut suiure la mesme regle pō detacher vt la, \ill{} la re, des chants de 6. notes qui font 720 varietez.
\note{« des 4320 chants » : lu tel quel ; la suite du calcul (5040, 10080, 20160, « oster de 40320 ») suppose 40320 chants, c'est-à-dire l'ordre de huit notes. Le copiste abrège « pour » en « pō » (p surmonté d'un tilde) et « premier » en « p\textsuperscript{r} ».}

p. 131. au commencem\textsuperscript{t}.

La combinaison des changem\textsuperscript{t}. est de 2 sortes, l'vne ou toutes les choses sont differentes co\textsuperscript{e} au\add{x} 12 cartes du picquet, l'autre ou elles peuuent estre indifferem\textsuperscript{t} ou toutes differentes ou toutes semblables, ou partie differentes, et partye semblables, co\textsuperscript{e} si l'on prenoit \uncertain{lesd.} ieux de 12 cartes en 12 ieux de picquet. Pour l'vne et l'autre il faut \uncertain{f\textsuperscript{e}} 12 nombres par multipli\add{cation} y compris le p\textsuperscript{r} qu'on multiplie, qui est 36, mais pō le p\textsuperscript{r} ou tout est different, il fau\add{t} multiplier par le nombre inferieur scauoir par 35, et le produit par 34 et cæt. Et pō le 2\textsuperscript{nd}. ordre, il faut multiplier par les superieurs, scauoir par 37, 38 et cæt. Donc pō le \ill{} il faut \uncertain{f\textsuperscript{e}} 11 multiplicateurs et le dernier produit \struck{donne} 599555620984320000 donn\add{e} la varieté des 12 cartes prises en 36 dans l'ordre, c'est a dire supposant qu'on les arr\ill{} en toutes les façons possibles, qui est le p\textsuperscript{r} cas, ou 1\textsuperscript{re} sorte de combinaisons mesl\ill{} qui suppose toutes les choses differentes, mais prises en vn plus grand nombre, et suppos\add{e} aussy l'ordre. partant pō auoir lesd. ieux de picquet sans l'ordre, puis qu'il ne change point le ieu, il faudra diuiser le grand nombre precedent qui a 5 quadril\ill{} pō son premier chefre, par la combinaison de l'ordre de 12 choses scauoir 479001\ill{} et on aura 1251677700 varietez de ieux de picquet. pō ce qui est de iouer et ietter l\add{es} cartes sur la table, il faut auoir egard a l'ordre, et chasque ieu se peut iouer en 479001600 façons qui est l'ordre de 12 et partant pō auoir en tout en combien \ill{} sortes on peut iouer les 12 cartes prises en 36, il faut multiplier lesd. 1251677700 \ill{} 479001600 pō auoir le grand nombre cydeuant trouué 599555 et cæt.
\note{« il faut f\textsuperscript{e} 12 nombres », « il faut f\textsuperscript{e} 11 multiplicateurs » : lu comme l'abréviation de « faire » (vue 22, « a f\textsuperscript{e}. ») ; la lettre initiale n'est pas sûre et pourrait être le p de « p\textsuperscript{r} » (« premier »). Le produit $36 \cdot 35 \cdots 25$ vaut bien 599555620984320000, et son quotient par $12! = 479001600$ vaut 1251677700.}

Voyla pō le 2\textsuperscript{me} changem\textsuperscript{t} ou toutes les choses sont differentes soit sans l'ordre ou av\add{ec} l'ordre. mais lors qu'elles peuuent estre semblables ou en tout ou en partye, il fau\add{t} multiplier \uncertain{37} par 37, et puis le produit par 38, iusques a ce qu'on ayt 12 nombres o\add{u} produits, autant que la multitude des choses, et ainsy le dernier nombre qui multipli\ill{} sera 47, et puis il faut diuiser le produit par l'ordre de 12. De mesme pō les 6 pani\add{ers} de fruits, desquels on en choisit 100, auec liberté de les prendre tous semblables et on v\ill{} mesme panier, ie prens 6 pō le terme des multipli. et le multiplie par 7 j'ay 42 puis 42 par 8 et cæt. tant que j'aye cent nombres y compris 6. le dernier nom\add{bre} multipliant sera 105, qui surpasse 6. de 99, scauoir de vn moins que la multi\add{tude} des choses 100, a cause que 6 est conté pō le p\textsuperscript{r}, le dernier produit estant diuisé \ill{} l'ordre des 100 choses, donnera le nombre requis.
\note{« multiplier 37 par 37 » : le premier chiffre est lu 37 ; le calcul décrit (douze facteurs jusqu'à 47) demande 36.}

Lors que les 12 cartes sont prises en 12 ieux pō estre semblables, et qu'on y garde l'ord\add{re} il se faut seruir des puissances quarrées, et sans l'ordre des puissances triangulai\add{res} les quarrées sont les ord\textsuperscript{res}: la triangulaire se fait par l'addition des puissances q\ill{} ont 1 moins d'exposant depuis la p\textsuperscript{re} qui est 1, iusques a celle qui a pareille rac\add{ine} co\textsuperscript{e} la 6\textsuperscript{me} puissance triangulaire de 5 est la so\textsuperscript{e} des cinq premieres cinquiesmes puissances, et la 5\textsuperscript{me} puissance de 5 est la so\textsuperscript{e} des cinq premieres quatriesmes puiss\add{ances} et la quatriesme puissance de 5 est la so\textsuperscript{e} des cinq premiers tetraedres, ou troisiesmes puissances, et le tetraedre de 5 est la so\textsuperscript{e} des cinq premiers triangles: co\textsuperscript{e} le 5 triangle ou le triangle de 5 est la so\textsuperscript{e} des cinq premiers nombres. voyez la tab\add{le} desd. puissances triangulaires iusques a la 12\textsuperscript{me} puissance de 25 pag. 136 de l'harmonie en latin, en la 12\textsuperscript{me} perpendiculaire, donc pō auoir les 12 cartes sans l\ill{}
\note{La phrase se poursuit en tête de la vue 22. Ici « puissances quarrées » désigne les puissances ordinaires et « puissances triangulaires » les nombres figurés.}

\page{22}
\note{Page de droite (recto), chiffrée « 11 » à l'encre en haut à droite. À gauche de la vue paraît le bord de la page précédente (vue 21), qui ne livre rien de plus que la vue 21 elle-même.}

Il faut prendre la 12\textsuperscript{me} puissance triangulaire de 36; qui est a la fin de lad. table ascauoir 1251677700. D'ou il est aisé de tirer la regle pō auoir les puissances triangulaires. par exemple la 6\textsuperscript{me} puissance triangulaire de 5 la Racine est 5, et l'exposant 6, prenez 6 nombres de suite dont la R. sera le moindre, scauoir est 5. 6. 7. 8. 9. 10. qu'il faut multiplier l'vn par l'autre, et diuiser le produit par l'ordre de la multitude desd. nombres qui est representé par l'exposant 6., et cest ordre est 720, ou pō euiter la diuision, ostez des 6 nombres susd. les six premiers nombres 1, 2, 3, 4, 5, 6, et il restera 7. 3. 10. dont le produit 210 est la puissance requise, scauoir la 6\textsuperscript{me} puissance triangulaire de 5 co\textsuperscript{e} on void a la table cotée de l'harmonie en latin et qui se void aussy dans ce liure pag. 145.
\note{Suite de la phrase de la vue 21 (« donc pō avoir les 12 cartes sans l'ordre, »). « R. » : la racine.}

p. 135 prop. 12\textsuperscript{me}.

Si l'on auoit 6 paniers pleins chacun d'vne espece de \struck{\ill{}} fruits differens, mais esgaux entre eux, co\textsuperscript{e} l'vn de poires, l'autre de prunes et cæt. toutes esgales, et qu'on voulut choisir vn cent desd. fruits dans lesd. paniers, dont chascun ayt du moins vn cent de fruits, afin qu'il se puisse si on veut prendre vn cent de mesme sorte, il faut f\textsuperscript{e} la combinaison generale de cent choses prises en 6 sortes de choses, car le 6 est R. c'est a dire la varieté des choses, et 100 est l'exposant, c'est a dire la multitude des choses, partant il faut prendre la 100\textsuperscript{me} puissance de 6. pō la varieté des \uncertain{diuerses} façons dont lesd. fruits pris en 6 panniers se peuuent arranger.

La combinaison generale contenant celle de l'ordre en toutes les façons possibles tant des choses differentes que des semblables, ne peut estre diuisée par led. ordre, parce que plus\textsuperscript{rs}. diuiseurs estans assemblez sont autre chose que diuisez: or l'ordre y est compris en toutes sortes de façons par exemple si on a 4 choses, elles y sont considerées confusem\textsuperscript{t}. soit qu'on les prenne toutes differentes co\textsuperscript{e} abcd, ou 2 differentes et 2 semblables co\textsuperscript{e} aabc, ou 3 semblables et vne autre co\textsuperscript{e} aaab, ou deux d'vne sorte et 2 d'vne autre co\textsuperscript{e} aabb, ou enfin toutes 4 semblables co\textsuperscript{e} aaaa, ou bbbb; Abcd se change en 24 façons, aabc en 12, aaab en 4, aabb en 6, et finalem\textsuperscript{t}. aaaa d'vne sorte; or tous ces nombres font 47, nombre p\textsuperscript{er} qui ne peut diuiser aucune puissance (dont la combinaison generale est faite) si elle n'auoit 47 pō racine. De plus l'ordre de 4 ayant 5 varietez, il faudroit diuiser 47 par 5, ce qui est trop incommode.

p. 138 ligne 21 apres ces mots (et la 5 du systeme

parce que le b fa $\natural$ mi fait deux degrez dans la la table, de sorte que ce fa est en C sol vt fa, qui est le 5\textsuperscript{me} degré ou la 5\textsuperscript{me} chorde de ce systeme
\note{« b fa $\natural$ mi » : entre « fa » et « mi », un signe en forme de 4, transcrit par le bécarre. « dans la la table » : le mot est répété.}

p. 139. ligne premiere apres ces mots (et la 15\textsuperscript{me} du systeme)

et partant il faut multiplier 484 qui est le 3\textsuperscript{me} nombre de la table par 15, ce qui donne 7260, et ainsy des autres: par exemple fa qui est la 4\textsuperscript{me} prise tousiours a rebours, est la 16\textsuperscript{me} du systeme, et partant il faut multiplier le 4\textsuperscript{me} nombre de la 5\textsuperscript{me} colomne de la table 10648, par 16. De mesme la 5\textsuperscript{me} note, mi qui est encore la 15\textsuperscript{me} du systeme, car il faut multiplier le nombre cinquiesme de la 5\textsuperscript{me} colomne a scauoir 234276, \add{par 15} dont il viendra 3514140, et non 3513840, co\textsuperscript{e} il y a dans le texte qui est fautif, c'est pourquoy il faut refaire la supputation.
\note{« par 15 » est écrit au-dessus de la ligne, après 234276, par le copiste ; l'\add{} le place. « co\textsuperscript{e} il y a dans le texte » : le copiste renvoie à une faute de l'imprimé.}

p. 144 au corollaire 1\textsuperscript{er}

\marginal{1\textsuperscript{er}} par les nombres on sout mille dificultez: comme lors qu'on propose 50 ou cent choses distinctes chacune d'vne toise, qu'il faut toutes porter dans vn mesme lieu, il faut seulem\textsuperscript{t}. multiplier le nombre des choses, par le nombre qui est moindre immediatem\textsuperscript{t} par exemple s'il y a 6 choses, estans multipliez par 5 l'on a 30 toises de chemin a f\textsuperscript{e}., et si le lieu ou il faut porter est esloigné de 7 toises, estans multipliez par 6, nous auons 42 toises, donc s'il y a 50 choses a ramasser, il faut multiplier 50 par 49, pō auoir 2450, et s'il y a cent, l'on aura 9900.

\marginal{2\textsuperscript{me}} Si l'on veut scauoir la so\textsuperscript{e} des nombres qui se suiuent immediatem\textsuperscript{t} co\textsuperscript{e} 1. 2. 3. 4. et cæt. si le dernier est pair co\textsuperscript{e} 14, l'on prend sa moitié 7 et puis on multiplie le nombre
\note{Les numéros « 1\textsuperscript{er} » et « 2\textsuperscript{me} » sont dans la marge de gauche, en regard des deux alinéas. La phrase se poursuit à la vue 23.}

\page{23}
\note{Page de gauche (verso) ; aucun chiffre en tête. Comme à la vue 21, les fins de ligne passent dans le pli de la reliure.}

plus grand de l'vnité, a scauoir 15 par 7 pō auoir 105. et si le dernier nombre \add{est} impair, comme 5 il faut le multiplier par son milieu 3, et l'on 15 de mesme. 15 doibt estre multiplié par 8 pō auoir 120 que font les nombre 1, 2, 3 et cæt iusques a 15.
\note{Suite du second corollaire de la vue 22. « et l'on 15 de mesme » : lu tel quel, un verbe semble manquer.}

p. 146 au commencem\textsuperscript{t}.

Il faut prendre garde que le 2\textsuperscript{me} terme apres 36, c'est a dire que le produit de 36 \add{par} 35 est 1260, mais ce nombre estant diuisé par les combinations ord\textsuperscript{res} de l'ordre il \ill{} 630: de sorte qu'il faut accommoder ce discours aux tables suiuantes autrem\textsuperscript{t} il seroit faux que le 2\textsuperscript{ond} terme de la multiplication fut 630, et le dernier 125677700, desq\ill{} les combinaons de la 3\textsuperscript{me} colomne sont ostez, ceux de la 4\textsuperscript{me} colomne sont les veritables produits, qui contiennent le changem\textsuperscript{t}. de l'ordre, tel qu'il peut estre seruant et iettant les chartes sur la table
\note{« 125677700 » : neuf chiffres, lus nettement ; la vue 22 donne pour le même nombre 1251677700.}

Ibid. prop. 17\textsuperscript{me}

Il y a icy vne combination a appliquer qui est lors qu'on prend par exemple 12 cha\ill{} du jeu de picquet dans 12 ieux de cartes differents ou dans le nombre des choses il y \add{a} tousiours a chasque fois quelqu'vne differente, co\textsuperscript{e} si l'on veut scauoir en comb\add{ien} de manieres on peut f\textsuperscript{e} vne compagnie de cent soldats, pris dans vn regim\textsuperscript{t} de mi\ill{} hommes. Pour les 12 cartes prises en 36, lors qu'on a douze ieux seulem\textsuperscript{t}, il faut prend\add{re} 36 pō racine, et 12 pō exposant et partant ce sera la 12\textsuperscript{me} puissance de 36, si on ne prenoit qu'vne carte on auroit que 36 varietez, si on prenoit 2 cartes on auroit 1296 varietez, scauoir le quarré de 36; pō trois cartes on prend le cube, pō quatre le quarr\ill{} et cæt. a l'infiny. Neantmoins autrepart m\textsuperscript{r}. Fren. \add{dit} que pō ces 12 cartes prises en \ill{} ieux de 36 cartes il faut multiplier 36 par 37, et le produit par 38 et cæt. iusques a \ill{} multiplications, ou le dernier multipliant sera 47: et que pō en oster l'ordre il faut diuiser par la combinaison ord\textsuperscript{re} de 12, de mesme que pō scauoir en combien de mani\ill{} on peut prendre 100 fruits en 6 paniers, il faut multiplier 6 par 7, et le produit \ill{} 8, et ainsy consequem\textsuperscript{t}. iusques a auoir 100 nombres y compris 6, et ainsy \struck{\ill{}} le dernier nombre multipliant est est 105, qui surpasse 6. de 99: et le dernier produ\add{it} diuisé par l'ordre de 100 choses, donne le nombre requis / Il faut donc concilier ces façons, et appliquer ce qu'il y a a dire/
\note{« dit » est écrit au-dessus de la ligne, sur « Fren. ». L'abréviation « Fren. » n'est pas développée sur la feuille ; y lire Frenicle est une inférence de cette lecture. « est est » : le mot est répété.}

\marginal{\uncertain{qa}} Ainsi 2 choses prises en 10, 2 a 2 font 100 varietez, car le nombre 100 contient 10 et 2, \ill{} 10 est sa racine et 2 son exposant. Et la varieté de 2 prise en \struck{\uncertain{neuf}} nombres, n'est qu\add{e} 81. si on ne prend qu'vne chose de 9, on n'aura que 9 varietez, si l'on en prend 2 cho\add{ses} puis qu'apres chascune des 9 choses on peut prendre chascune des mesmes 9 l'vne ap\add{res} l'autre il faut pō auoir ceste varieté multiplier 9 par 9 pō auoir 81: et si on pr\ill{} 3 choses, on pourra deuant chascune des 81 lettres precedente mettre chascune \ill{} 9 choses, partant il faut multiplier 81 par 9 pō auoir 729 cube de 9 et ainsy \ill{} reste; car pō 4 choses il faut prendre le quarrequarré, c'est a dire multiplier 729 \ill{} 9 et cæt. Notez que la diuersité des choses qu'on prend, co\textsuperscript{e} icy 9 est tousiours racine, et la multitude des choses sert d'exposant.
\note{Dans la marge de gauche, en regard de la ligne qui précède « Ainsi », deux lettres, « qa » ou « ga », d'une encre semblable ; sens non établi.}

p. 152 prop. 20.

Si l'on peut faire que les 7 notes differentes de chasque partye puissent tousiours fa\ill{} des accords ensemble, ou des discords supportables suiuant les loix de la musique il \add{y} aura 2520473760000, d'ont l'on n'en pourra chanter que 45048300 dans le te\ill{} de cent ans, bien que chasque chant ne dure qu'vne seconde minute, de sorte qu'il faudroit plus de cent mille siecles pō chanter tous ces chants iour et nuit faisant 4 chants pō accorder tousiours ensemble, ilz seruiroient aussy pō long te\ill{} car l'exemple de 4 notes qui font trois parties contient vne 26\textsuperscript{me} de celuy, \ill{} 4 partyes vne 36\textsuperscript{me} ou 5 octaues. Les notes sont les cadences du p\textsuperscript{r} mode, ou d'vn \ill{} par exemple, vt, mi, sol, fa, iointes aux 4 notes d'vne taille differentes des 4 notes precedentes font 576 duo. les 4 de la hautecontre font vn trio de 13729 chants et le dessus fait 319376 quatuor. Neantmoins vne grande partye des accords de ces varietez ne valent rien a cause des quartes, des 5 et des 8 qui se suiuent souuent 2 \ill{} suite, et des interualles desagreables et dificiles a chanter, par lesquels on passe d'vne consona\ill{} a vne autre soit a 2, a 3, ou a 4. Il n'y a point d'autres notes autentiques au nombre de 4, qu\add{i} puissent tous\textsuperscript{rs} accorder a 2, 3, et 4, que les 2, 3, ou 4 cadences des modes, et si chasque note dure u\ill{} mesure, l'on aura pō plus 12 \ill{} a chanter./
\note{« 13729 », « 319376 » : lus tels quels ; ils ne sont pas les puissances de 24 (13824, 331776) auxquelles 576 = $24^2$ ferait attendre.}

\note{Au pied de la page, à droite, la réclame « pag. 154 au droit », qui est la rubrique de tête de la vue 24.}

\page{24}
\note{Page de droite (recto), chiffrée « 12 » à l'encre en haut à droite. Au pied de la page, à droite, un « 2 » isolé, de la même encre, peut-être une marque de cahier.}

pag. 154 au droit de la musique

parce qu'il n'y a d'ordinaire dans les beaux traits de chants que 3 ou 4 notes, celles de la 6. de la 9\textsuperscript{me} propos. qui se varient en 720 manieres Et les 8 du gros volume qui se varient en \struck{40}40320 manieres, peuuent seruir pō f\textsuperscript{e} de beaux chants, quoy qu'ilz dependent souuent de la repetition des mesmes notes./

\marginal{\ill{} 1\textsuperscript{er} des \ill{}sonances}
\note{Dans la marge de gauche, sous un trait horizontal qui sépare cette remarque de la précédente, un titre courant dont le début passe dans le pli : « … 1\textsuperscript{er} des …sonances », sans doute « Livre 1\textsuperscript{er} des Consonances ». Les renvois qui suivent (p. 26, 46, 53…) repartent donc d'une nouvelle pagination.}

p. 26. prop. 6. lig. 14.

Ce qui se peut encore expliquer par le souflem\textsuperscript{t} duquel on vse pō esbranler vne lampe d'eglise, vne cloche, ou autre chose suspendue, car le 2\textsuperscript{me} soufle renforce le p\textsuperscript{r} et cæt. comme explique Galilée en son 1\textsuperscript{er} dialogue.

p. 46. coroll. 1\textsuperscript{er}

Galilée remarque dans son 1\textsuperscript{er} dialogue qu'il a veu des petites ondes sur vne lame de cuiure qu'il racloit d'vn ciseau, et que ces petites ondes paralleles faisoient des sons plus aigus en mesme raison que leur \uncertain{estresse} les rendoit plus proches les vnes des autres, ce qui reuient aussy au verre frotté sur le bord, lequel sautant a l'octaue les ondes de l'eau qui fremit se fendent en deux, il dit encore que le son de la lame faisoit trembler des chordes d'espinette a la quinte.

Voyez les deux dernieres pages de la table des matieres, ou les chordes pendues a des clous representent les consonances par leur mouuem\textsuperscript{t}.

p. 53. lig. 28.

Kircker p. 831 de magnetismo musica, explique l'vnion des retours dans l'octaue, la quinte, la quarte et l'vnisson, apres Galilée dans ses liures du mouuement.

p. 55. au commencement

l'octaue ou la raison de 2 a 1, ne se peut diuiser en trois surparticuliers qu'en cinq sortes co\textsuperscript{e} l'on peut voir a la marge
\[
\begin{array}{|c|c|c|}
\hline
\cdot 3 & 4\cdot 5 & 15\cdot 16 \\ \hline
\cdot 3 & 5\cdot 6 & 9\cdot 10 \\ \hline
\cdot 3 & 6\cdot 7 & 7\cdot 8 \\ \hline
\cdot 4 & 3\cdot 4 & 8\cdot 9 \\ \hline
\cdot 4 & 4\cdot 5 & 5\cdot 6 \\ \hline
\end{array}
\]
\note{Table dans la marge de gauche, en regard de la remarque, trois colonnes réglées. Le premier chiffre de chaque rapport de la première colonne passe dans le pli ; on ne lit que « ·3 » trois fois et « ·4 » deux fois.}

p. 58. prop. 15.

Notez que ceste proposition se doibt entendre des vrais interualles, suiuant n\textsuperscript{re} theorie, autrem\textsuperscript{t}. elle seroit fausse, car si l'on suppose les tons egaux, co\textsuperscript{e} sont les instrumens ord\textsuperscript{res}, la quinte triplée seroit la 13\textsuperscript{me}, ou sept\textsuperscript{e} maieure repetée, co\textsuperscript{e} la tierce trois fois repetée feroit l'octaue et cæt. ce qui ne peut arriuer en la vraye theorie ou les tons sont inegaux, aussy bien que les demytons: ce qu'il faut soigneusem\textsuperscript{t} remarquer contre les praticiens

lig. 6. deuant la fin.

Notez qu'vne chorde ou vn tout estant diuisé en deux moitiez, si l'vne des moitiez est diuisée en 2, et puis en 4, 8, 16 et cæt. a l'infiny, toutes ces parties seront esgales a toutes les chordes: et si le tiers est diuisé en tous ses autres tiers, tous ses tiers a l'infiny feront la moitié du tout, co\textsuperscript{e} tous les quarts $\frac{1}{3}$, tous les dixiesmes $\frac{1}{9}$ a l'infiny

p. 64. lig. 4. auant la fin

l'on a mainten\textsuperscript{t}. des pratiques qui font paroistre la quarte bonne en duo, voyez la fantaisie du 5\textsuperscript{me} liure de la Composition page 200.

p. 71. lig. 9.

tromperie notable dificile a euiter

lig. 22

Il vaut mieux dire que le p\textsuperscript{r} coup de la grosse chorde ne s'vnit pas auec le p\textsuperscript{r} de la deliée mais seulem\textsuperscript{t}. son 3\textsuperscript{me} coup auec le 4\textsuperscript{me}, et partant la deliée a trois coups qui ne s'vnissent point, et trois fois plus de coups desunis que d'vnis, co\textsuperscript{e} la quinte en a deux fois plus de desunis, et l'octaue en vnit la moitié, en ayant autant d'vnis que de desunis, il faut dire la mesme chose des autres consonances a proportion

p. 72. l. 11.

Il est mieux dit que la quinte desunit deux de ses mouuem\textsuperscript{ts} du son aigu, et la quarte trois, et partant que la quinte est sesquialtere de la quarte en agrem\textsuperscript{t}, et sous double de l'octaue. Tout ce qui a fait faillir la speculation est de supposer que les premiers coups de chasque corde s'vnissent ensemble. ce qui suit depuis mais et cæt. est meilleur

l. 30.

Ce qui empesche que le premier coup de la grosse chorde ne s'vnit pas auec le p\textsuperscript{r} de la deliée,
\note{La phrase se poursuit à la vue 25.}

\page{25}
\note{Page de gauche (verso) ; aucun chiffre en tête. Les fins de ligne passent dans le pli.}

est qu'elle n'acheue pas son premier tour auec le premier tour de la sienne
\note{Fin de la remarque « l. 30. » de la vue 24.}

p. 74 prop. xxvij\textsuperscript{me}

voyez les 27 diuisions que la quarte a en trois raisons surparticulieres sans en pouuoir souffrir dauantage page 142 du 3\textsuperscript{me} liure des genres

p. 76. prop. xxx\textsuperscript{me}

Il vaut mieux dire que la tierce n'est surmontée par la quarte que d'vn quart, ou de raison de 5 a 4 puis que la tierce n'vnist qu'\struck{\ill{}} \uncertain{$\frac{1}{5}$} de ses temps, et la quarte le quar\add{t} des siens

p. 91. au commencem\textsuperscript{t}.

\uncertain{L'autheur} traicte amplem\textsuperscript{t} de ces trois mediettez dans son 3\textsuperscript{me} liure depuis la 5\textsuperscript{me} propo\textsuperscript{n} iusq\ill{} a la 27\textsuperscript{me} ou il donne le moyen de les trouuer dans le demicercle, entre 5 lignes et entre 5 nombres, dont les moindres en double raison sont 12. 9. 6. 4. 3. et en raison triple 18. 12. \ill{} car 12, 9, 6: ou 18, 12, 6: sont en proportion arithmetique, car 12 est le milieu, et 6. est m\ill{} geometrique entre 9 et 4 comme 4 est milieu harmonic entre 6 et 3 et separem\textsuperscript{t} les moindres termes de la mediété arithmetique sont 3, 2, 1, dans la geometrique 4, 2, 1 et \ill{} l'harmonique 6, 4, 3: Mais il explique 7 autres medietez, et donne la maniere de le\add{s} trouuer par l'analogie geometrique laquelle seule merite ce nom.

p. 92. l. 13.

Je rends ceste 5\textsuperscript{me} maniere generale en vn au\textsuperscript{re} lieu, ascauoir dans la 5\textsuperscript{me} remarque \ill{} la preface mise au commencem\textsuperscript{t}. de ce liure, ou j'explique 4 manieres pō trouuer milieu harmonic\ill{}

p. 109. au commencem\textsuperscript{t}.

Table 57. La regle de trois sert pō trouuer les termes de chasque interualle, par exemple si on veut scauoir quel nombre fait le ton maieur auec 25, on trouue $22\frac{2}{9}$, en disant \ill{} 9 donne 8 combien 25. Mais si l'on veut que ce nombre soit quarré co\textsuperscript{e} 25, de mesme que 81 et \ill{} qui sont les quarrez de 9 et 8 sont en raison double de 9 a 8 il prend la racine de tel nombre quarré qu'on veut par exemple 2 racine de 4, mais parce que 4 ne peut auoir vn nombre sesquioctaue sans fraction, il prend les quarrez des deux nombres originaux \ill{} ascauoir 81 et 64 et dit si 81 donne 64 que donneront 4, on aura $3\frac{13}{81}$ qui est nombre duquel la racine \uncertain{est} sesquioctaue de 2 et la R. de $3\frac{13}{81}$ se trouue en disant si 9 (R. de 81) \ill{} 8 (racine de 64) combien donnera 2 (R. de 4) on a $1\frac{7}{9}$, lequel multiplié par soy mesme fait $3\frac{13}{81}$ partant 2 est sesquioctaue de $1\frac{7}{9}$. L'on peut plus briefuem\textsuperscript{t}. trouuer ce nombre en disant si 9 donne 8 combien 2./
\note{« Je rends » : le copiste reproduit la première personne de Mersenne, ici et plus loin (« ou j'explique »).}

\marginal{Du liure 3\textsuperscript{me} des consonnances}
\note{Titre courant dans la marge de gauche, sous un trait horizontal terminé en flèche vers la gauche.}

p. 142 au commencem\textsuperscript{t}.

La quarte ou la raison de 4 a 3 se peut diuiser 27 fois en trois raisons surparticulieres bien que Briennius n'en tient que 15 possibles. voicy lesd. 27 manieres qui font 27 differentes especes de tetrachorde
\[
\begin{array}{ccc}
\underset{1}{9\cdot 10/10\cdot 11/11\cdot 12} & \underset{2}{9\cdot 10/8\cdot 9/15\cdot 16} & \underset{3}{4\cdot 5/16\cdot 17/255\cdot 256} \\[4pt]
\underset{4}{4\cdot 5/17\cdot 18/135\cdot 136} & \underset{5}{4\cdot 5/18\cdot 19/\ldots} & \underset{6}{4\cdot 5/19\cdot 20/75\cdot 76} \\[4pt]
\underset{7}{4\cdot 5/20\cdot 21/63\cdot 64} & \underset{8}{4\cdot 5/21\cdot 22/55\cdot 56} & \underset{9}{4\cdot 5/23\cdot 24/45\cdot 46} \\[4pt]
\underset{10}{4\cdot 5/25\cdot 26/3\ldots} & \underset{11}{4\cdot 5/27\cdot 28/35\cdot 36} & \underset{12}{4\cdot 5/30\cdot 31/31\cdot 32} \\[4pt]
\underset{13}{15\cdot 16/5\cdot 6/24\cdot 25} & \underset{14}{15\cdot 16/6\cdot 7/14\cdot 15} & \underset{15}{5\cdot 6/10\cdot 11/99\cdot 10\ldots} \\[4pt]
\underset{16}{5\cdot 6/11\cdot 12/54\cdot 55} & \underset{17}{5\cdot 6/12\cdot 13/39\cdot 40} & \underset{18}{5\cdot 6/14\cdot 15/27\cdot 28} \\[4pt]
\underset{19}{5\cdot 6/18\cdot 19/19\cdot 20} & \underset{20}{9\cdot 10/6\cdot 7/35\cdot 36} & \underset{21}{9\cdot 10/7\cdot 8/20\cdot 21} \\[4pt]
\underset{22}{6\cdot 7/8\cdot 9/63\cdot 64} & \underset{23}{6\cdot 7/11\cdot 12/21\cdot 22} & \underset{24}{6\cdot 7/14\cdot 15/15\cdot 16} \\[4pt]
\underset{25}{7\cdot 8/7\cdot 8/48\cdot 49} & \underset{26}{7\cdot 8/8\cdot 9/27\cdot 28} & \underset{27}{7\cdot 8/12\cdot 13/13\cdot 14}
\end{array}
\]
\note{Tableau réglé, cinq divisions par rangée (ici trois), chacune numérotée sous ses trois raisons ; les divisions sont séparées sur la feuille par une double barre. Les derniers chiffres des divisions 5, 10 et 15 passent dans le pli. Division 23 : le « 11 » est repassé sur un autre chiffre.}

p. 144 au commencement

\note{Un schéma du grand système : sur une ligne, treize cases séparées par des traits, chacune donnant la lettre et la syllabe de la note suivies de deux lettres grecques abrégées (sans doute l'initiale du nom grec de la corde), que cette lecture ne déchiffre pas : « a re \ill{} | G sol \ill{} | f fa \ill{} | E mi \ill{} | D sol \ill{} | C fa \ill{} | $\natural$ mi \ill{} | a re \ill{} | G sol \ill{} | f fa \ill{} | E mi \ill{} | D sol \ill{} | C fa \ill{} », la ligne se poursuivant sur la page suivante. Des accolades tracées sous la ligne réunissent les cases en tétracordes, légendés « des derniers » (écrit d'abord, puis biffé avec « ton », à l'extrême gauche), « ton », « des disioincts », « des conioincts » (celui-ci placé sous une rangée supplémentaire « D \ill{} | C \ill{} | b \ill{} »), « ton », « des moyens », « des basses ». Deux mots sont écrits verticalement : à gauche, \ill{} ; au milieu, à la jonction des tétracordes, « \uncertain{Moyenne} ». À droite, au bord du pli, « A\ill{} ».}

Les premieres et les dernieres chordes des tetrachordes doibvent tous\textsuperscript{rs} demeurer immobiles
\note{Au pied de la page, à droite, la réclame « a scavo\ill{} » ; la phrase se poursuit à la vue 26.}

\page{26}
\note{Page de droite (recto), chiffrée « 13 » à l'encre en haut à droite. Des traits et des notes vus par transparence (le verso) courent sous l'écriture.}

a scauoir les mi ou les re ou la: mais la seconde ē mobile dans l'enharmonique seulem\textsuperscript{t}., parce qu'il faut l'abbaisser d'vne diese, et le sol ē mobile au chromatiq, ou il s'abbaisse d'vn demyton, et dans l'enharmonic d'vn ton entier, de sorte qu'il est vnisson du fa diatoniq et chromatiq, au lieu que mainten\textsuperscript{t} pō f\textsuperscript{e} le sol chromatiq on hausse le fa d'vn demy ton.
\note{Suite de la remarque « p. 144 » de la vue 25 (« doibvent tous\textsuperscript{rs} demeurer immobiles »). « ē » : le copiste abrège « est » par un petit signe bouclé posé sur la ligne.}

p. 145 au commencem\textsuperscript{t}.

voyez les deux premieres propo\textsuperscript{ns} du 6\textsuperscript{e} liure de la Composition, ou il y a de nouuelles eschelles de musique, et ou celle cy ē expliquée, au lieu des six syllabes de guy Aretin, on en met mainten\textsuperscript{t} 7. pō faire les muances, et pō ce suiet le maire fait sa \uncertain{taradirie} qu'il nomme de ces mots, ta, ra, ma, fa, sa, la, za, dont il diese ta, ra, fa, za, afin de f\textsuperscript{e} le demiton partout co\textsuperscript{e} tous les instrumens, il prend tous\textsuperscript{rs} le ta, ou vt en f fa vt, de sorte que le za est vn ton sur le la par b mol, et pō lors il ne faut pas dieser le fa, mais bien le za, et lors qu'on chante par $\natural$ quarré on diese le fa et le za voicy ses caracteres \ill{} et auec valeur de notes \ill{}, qu'il ponctue ou virgule differemm\textsuperscript{t} pō les faire finir a l'estendue de 4 ou 5 octaues, et chascun confesse que la composition ē plus aisée auec ses caracteres a raison que n'ayant qu'vne seule ligne pō chasque partye, l'on voit plus plus viste et plus aysem\textsuperscript{t}. les consonances ou dissonances que l'on fait, mais j'ay donné ma methode de composer en ceste façon dans les 2, 3, et 4 propo\textsuperscript{ns}. du 6\textsuperscript{e} liure de la composition \struck{qu'il} \add{et} vaut mieux mettre le ba, bi d'\uncertain{olinde}./
\note{Les « caracteres » sont dessinés dans la ligne : d'abord sept signes simples, ressemblant à V, T, U, O, S, $\sharp$, Z ; puis les mêmes signes « avec valeur de notes », munis de hampes et de crochets. Ils ne sont pas reproduits. « et » est écrit au-dessus de « qu'il » biffé. « plus plus » : le mot est répété.}

p. 146. au commencement

\marginal{A} Jean le Fevre secret\textsuperscript{re} du Cardinal de Givry euesque de Langres a tous\textsuperscript{rs}. suiuy dans sa triple musique escritte l'an 1563 les eschelles de Gafurus voicy les nombres dont il vse pō representer les six voix vt re mi fa sol la.

\marginal{B.} Les p\textsuperscript{ers} nombres representent les lignes du monochorde, le second les surfaces co\textsuperscript{e} quarrez des premiers, et les troisiesmes representent les pesanteurs co\textsuperscript{e} cubes des premiers. pō le demy ton il est de 256 a 243 et pō vser des termes de Gafure il y a adiousté les demy tons partout \& commence par a re

\marginal{\ill{}} Tous ces poids font en grains $86681\frac{17}{27}$ ou 9 liures 6 onces, 3 gros, 65 grains et $\frac{17}{27}$

\marginal{C} \struck{par} en la page 177 il vse encore de ces termes lineaires
\note{Les lettres A, B, C sont dans la marge de gauche et renvoient aux trois tables qui suivent ; devant « Tous ces poids », un signe de renvoi en forme de croix barrée.}

Table A
\[
\begin{array}{c|c|c|c|c}
\text{la} & 16 & E & 256 & 4096 \\
\text{sol} & 18 & D & 324 & 5832 \\
\text{fa} & 20\frac{1}{4} & C & 410\frac{1}{16} & 8303\frac{49}{64} \\
\text{mi} & 21\frac{1}{3} & \natural & 455\frac{1}{9} & 9709\frac{1}{27} \\
\text{re} & 24 & A & 576 & 13824 \\
\text{vt} & 27 & G & 729 & 19683
\end{array}
\]

Table B
\[
\begin{array}{c|l}
A & 4608 \text{ en grains } 8 \text{ onces} \\
\text{—} & 4854\frac{14}{27}\ \ 8 \text{ onc. } 3 \text{ gros } 30 \text{ grains} \\
G & 5184\ \ 9 \text{ onc.} \\
\text{—} & 5461\frac{1}{8}\ \ 9 \text{ onc. } 3 \text{ gro. } 6 \text{ gra.} \\
F & 5832\ \ 10 \text{ onc. } 1 \text{ gro.} \\
E & 6144\ \ 10 \text{ onc. } 5 \text{ gro. } 24 \text{ gra.} \\
\text{—} & 6472\ \ 11 \text{ onc. } 13\cdot 64 \\
D & 6912\ \ 12 \text{ onc.} \\
\text{—} & 7281\frac{1}{9}\ \ 12 \text{ onc. } 5 \text{ gro. } 9 \text{ gra.} \\
C & 7776\ \ 13 \text{ onc. et demye} \\
B & 8192\ \ 14 \text{ onc. } 1 \text{ gro. } 56 \text{ gra.} \\
\text{—} & 8748\ \ 15 \text{ onc. } 1 \text{ gro. et demy} \\
A & 9216\ \ 16 \text{ onc.} \\
\text{—} & 9709\frac{1}{27} \\
G & 10368
\end{array}
\]

Table C
\[
\begin{array}{c|c|c|c|l}
F & \text{fa} & 10\frac{1}{8} & A & 81 \\
G & \text{s} & 9 & G & 64 \\
A & \text{la} & 8 & F & 50\frac{46}{81} \\
\natural & \text{bi} & 7\frac{1}{9} & E & 45\frac{9}{16} \\
C & \text{u} & 6\frac{3}{4} & D & 36 \\
D & \text{r} & 6 & C & 28\frac{4}{9} \\
E & \text{m} & 5\frac{1}{3} & B & 25\frac{5}{8}\quad \frac{1}{1256} \\
F & \text{f} & 5\frac{1}{16} & A & 20\frac{1}{4} \\
G & \text{s} & 4\frac{1}{2} & G & 16 \\
A & \text{l} & 4 & &
\end{array}
\]
\note{Les trois tables sont côte à côte sur la feuille, séparées par des doubles traits verticaux. Dans la table B, le signe de l'once est rendu « onc. » ; les tirets de la première colonne marquent les degrés intermédiaires, sans lettre. Ligne E : un « 10 » plus gras est repassé sur un autre chiffre. Ligne « 6472 » : « 13·64 » est lu tel quel, sans unité. Table C : la fraction $\frac{1}{1256}$ est écrite à l'écart, à droite de la ligne B, et reliée à « $25\frac{5}{8}$ » par un petit trait ; le « 1256 » est douteux. « 9709 », dans la table A, porte sa fraction $\frac{1}{27}$ comme dans la table B.}

\note{Dans la marge de gauche, en regard de la remarque suivante, une petite portée porte une suite de notes rondes, dont deux précédées d'un dièse : l'exemple de la « quarte » dont il est question.}

p. 151. l. 7.

Dans ceste quarte la 1\textsuperscript{re} Et la 2\textsuperscript{de} chorde ē diatonique, la 3\textsuperscript{me} pure chromatique, la 4\textsuperscript{me} diatonique propre, la 4\textsuperscript{me} diatonique, la 5\textsuperscript{e} luy est aussy propre mais apartient non au ton ou au mode dorien, mais au phrygien selon m\textsuperscript{r}. Doni, et la 6, aussy bien que la 1\textsuperscript{re} et 2 sont communes au genre diatonique, et au chromatiq, la 5\textsuperscript{me} est le $\natural$ mi du phrygien

Selon l'ancien systeme d'Aretin et de pythagore la tierce maieure composée de 2 tons maieurs est de 81 a 64 la mineure de 32 a 27 et la sept\textsuperscript{e} maieure de 27 a 16.

\page{27}
\note{Page de gauche (verso) ; aucun chiffre en tête. La page porte quatre exemples notés, tracés à la plume par le copiste ; ils sont décrits, non reproduits.}

p. 154. a la fin

diuerse maniere d'entonner les notes du tetracorde diatonic, chromatic, enharmo\ill{} la 1\textsuperscript{re} hausse les chordes, et la 2\textsuperscript{me} les abbaisse: et qui est plus pour la lyre des anciens
\note{Une portée de cinq lignes, clé de fa, divisée en six mesures de quatre notes rondes (dièses et croix devant certaines), avec sous chaque mesure les syllabes : « mi fa sol la », « mi fa fa la », « mi fa fa la », « mi fa sol mi » (surmonté de quatre abréviations grecques, les mêmes qu'à la vue 25), « mi fa sol mi », « mi fa sol mi ».}

p. 171 lig. 15.
\note{Une portée, clé de fa, portant une gamme montante de notes rondes et noires, avec des dièses.}

Les notes noires signifient les feintes ou touches noires de l'orgue et de l'espine\ill{}

p. 180. prop. 15\textsuperscript{me}

on peut oster toute ceste propos\textsuperscript{on} 15 depuis la 16\textsuperscript{me} ligne, parce que l'on n'vse pas ordinairem\textsuperscript{t} de ces especes d'octaue

p. 181. ligne 6.

Il n'y aura que 21 especes, parce que la 7\textsuperscript{me} de la page precedente, et la 3\textsuperscript{me} de celle\ill{} sont semblables.

p. 182. au commencem\textsuperscript{t}

Au 6\textsuperscript{me} liure latin prop. 22, j'ay donné trois raisons pō lesquelles les modes commencent plustost par vt de C fa vt que par les au\textsuperscript{res} notes, a quoy vne 4\textsuperscript{me} raison ē a adiouste\ill{} a scauoir que dans ceste espece d'octaue, les 4 especes de quinte et les 3 de quarte \ill{} trouuent toutes immediatem\textsuperscript{t} les vnes apres les autres sans interruption, ce qui n'arriue a nulle autre espece d'octaue.

De plus il se trouue deux fois chasque espece de quarte dans ceste espece d'octaue.

Pour les quintes il ne s'y en trouue que deux especes de suite car la 3\textsuperscript{me}. espece manq\ill{} de repetition en haut donc selon le s\textsuperscript{r}. Vincent voicy le p\textsuperscript{r} et plus excellent mode
\note{Une portée, clé de fa, divisée par des barres, portant en accords de deux à trois notes rondes la suite des espèces du mode.}

Ce p\textsuperscript{r} mode est nommé lydien par quelques vns, voyez p. 186.

L'on void l'opinion qu'a le s\textsuperscript{r}. Dony des modes dans la 30\textsuperscript{me}. prop. du liure des cloches, \ill{} dans la 12 du 5\textsuperscript{e} liure de la Composition

De ces douze modes il y en a de semblables, car le p\textsuperscript{r} et le 8; le 2\textsuperscript{me} et le 9\textsuperscript{me} sont mesme ch\ill{} co\textsuperscript{e} le 3\textsuperscript{me} et le 10\textsuperscript{me}, le 4\textsuperscript{me} et le 11\textsuperscript{me}, le 5\textsuperscript{me} et le 12\textsuperscript{me}; de sorte qu'il n'y a que 7 mo\ill{} differents, suiuant les 7 lieux differents des demitons, co\textsuperscript{e} je remarque dans la pa\ill{} 188.

p. \struck{\ill{}} 183 apres la ligne 17.

La quatre espece d'octaue ne donne que le 7\textsuperscript{me} mode; et la 7\textsuperscript{me} espece ne donne que le 6 voyez la 12\textsuperscript{me} propo\textsuperscript{on} du 5\textsuperscript{me} liure de la composition.

Le s\textsuperscript{r} Dony commenceant le mode dorien en mi, dit qu'il respond a la musique des grecs hesychast\ill{} paisible et temperée, le mode phrygien commenceant par re a l'ecstatique ou enthousiast\ill{} et impetueuse. Et le lydien commenceant par vt a l'alegre diastaltique, comme le myxolydien systaltique, ou melancholique, ou languissant, suiuant ce que j'ay remarqué dans la 32\textsuperscript{me} propo\textsuperscript{on} du liure des instrumens de percussion.

p. 186 et 187 connexion modale des 7 tons du s\textsuperscript{r}. Dony
\note{Sept portées côte à côte, séparées par des traits, intitulées « Hypodorien », « Hypophrygien », « Hypolidien », « Dorien », « Phrygien », « Lydien », « Myxolydi\ill{} » ; chacune porte une clé de fa et une échelle montante de notes rondes et carrées. Sous chacune des six dernières, une accolade ramène la note initiale à celle de l'hypodorien, et le mot « hypodorien » est écrit verticalement en dessous.}

si chasque mode, ou harmonie, ou espece d'octaue, commencoit au mesme ton de voix que l'hypodorien, nous n'a\ill{} qu'vn mesme ton, hypodorien, mais parce que chasque espece d'octaue commence de plus haut en plus haut, le\ill{}
\note{Au pied de la page, à droite, la réclame « sont aussy » ; la phrase se poursuit à la vue 28.}

\page{28}
\note{Page de droite (recto), chiffrée « 14 » à l'encre en haut à droite.}

sont aussy bien icy variez que les 7 modes, ou harmonies, et si l'on commencoit de plus haut en plus haut sans varier les interualles de l'hypodorien, l'on chanteroit les sept tons differents dans vn mesme mode
\note{Fin de la remarque « p. 186 et 187 » de la vue 27.}

p. 189. lig. 11.

D'autres ayment mieux mettre le demiton mineur de D a d, le maieur de d a E le mineur de G a g le maieur de g a A et le moyen de B a $\natural$, et dans la seconde table de G a A le ton mineur d'A a $\natural$ le maieur ce qu'on verra depuis la 4\textsuperscript{me} propo\textsuperscript{on} iusques a la 13\textsuperscript{me}

p. 192
\[
\begin{array}{|c|c||c|c|}
\text{b mol} & & \natural\ \text{B quarre} & \\ \hline
\text{ta} & \text{fa} & \text{fa} & \text{vt} \\
\text{za} & \text{mi} & \text{mi} & \text{mi} \\
\text{la} & \text{re} & \text{re} & \text{la} \\
\text{sa} & \text{sol} & \text{fa} & \text{sol} \\
\text{fa} & \text{fa} & \text{mi} & \text{fa}\sharp \\
\text{ma} & \text{mi} & \text{re} & \text{mi} \\
\text{ra} & \text{re} & \text{sol} & \text{re} \\
\text{ta} & \text{vt} & \text{fa} & \text{vt} \\ \hline
\end{array}
\]
\note{Table dans la marge de gauche, en regard de « p. 192 », sans autre texte. À sa droite, deux mots douteux, « le $\natural$ » et « b l'\ill{} », en regard des deux premières lignes, et « fe » en regard de la ligne du fa dièse.}

p. 193. l. 10

Les dictions estans trop rudes l'on en peut inuenter d'autres plus douces

p. 194. coroll. 2\textsuperscript{me} lig. 6.

le \uncertain{9\textsuperscript{me}} a l'vne de ses cadences en d la re sol et non le second. Il faut donc icy prendre garde

p. 242. lig. 34

voyez ceste table au 5\textsuperscript{me} liure prop. 6. pag. 308.

p. 243. prop. 15 article troisiesme

par exemple les notes a queüe monstrent icy les suppositions des notes absentes
\note{Dans la marge de gauche, en regard, une petite portée de deux mesures porte quelques notes rondes, dont certaines à queue.}

p. 251 lig. 14

Ceste methode requiert que l'on aye tous\textsuperscript{rs}. la derniere, ou plus basse note de la basse dans l'esprit puis qu'il n'y a nulle note en chasque partye qui ne se raporte a elle

p. 255. prop. 20 ligne 4.

Ce qui se doibt entendre sous le signe du temps imparfait C, auquel se rapportent tous les autres signes tant d'augmenta\textsuperscript{on} que de diminution, parce que la maxime, aussy bien que les autres notes, changent de valeur a chasque signe different, mais j'ay seulem\textsuperscript{t}. icy mis le precedent, parce qu'il est le plus vsité de tous.

lig. 3\textsuperscript{me} deuant la fin

toutesfois l'on ne donne pas pō l'ord\textsuperscript{re}. vn signe par\textsuperscript{r} de pause a la maxime, mais on y met deux pauses de la longue

p. 295 lig. 25

Je l'apelle 3\textsuperscript{me} partye, parce que la 3\textsuperscript{me} minime du dessus du 4\textsuperscript{me} exemple bat contre les deux p\textsuperscript{ers} quarts, ou la moitye de la 2\textsuperscript{me} semibreue de la basse, de sorte qu'il reste la 3\textsuperscript{me} p\textsuperscript{tie} pō le point, ou le 3\textsuperscript{me} quart, et la 4\textsuperscript{me} p\textsuperscript{tie} ou le dernier quart ē pō batre contre la noire du dessus

p. 315 lig. 23 de la propo\textsuperscript{on} 8\textsuperscript{me}

\note{Un exemple à deux portées, en rondes, blanches et carrées, divisé par des barres, avec sous la portée supérieure les mots « cadences lydiennes » ; une accolade à droite réunit les deux portées, et la ligne qui suit est écrite en regard de l'accolade.}

Le tout selon m\textsuperscript{r}. Dony

\marginal{Des Instruments \ill{}chorde}
\note{Titre courant dans la marge de gauche ; la fin, « …chorde », est en partie dans le pli. Les renvois qui suivent (p. 3, p. 5…) repartent d'une nouvelle pagination.}

p. 3. propo\textsuperscript{on}. 2. ligne 5.

quelques vns co\textsuperscript{e} le sieur p. 151 de sa triste musique, tiennent qu'on a premierem\textsuperscript{t} fait les chordes de nerfs. et les chordes de boyau vallent mieux, quand les moutons sont tuez en pleine lune, et le pli qu'on fait en les pliant, fait quelque espece de dissolution en corrompant sa rondeur. estant vieille il s'y peut engendrer vne tigne, ou vermisseau qui la ronge, et lors elle sonne cas au lieu de la mangeure ou morsure; et puis si elle \uncertain{rauste} et \uncertain{chauste} elle est fausse. Il ne faut aussy laisser aux boyaux aucune particule de graisse ou de chair superflue, car cela \ill{} la chorde fausse. De mesme il croit que les cordes d'intestins de loup et de brebis sur mesme instrument ne s'accordent iamais, ce que nul n'ayant experimenté, \ill{} le mette entre les fables, puis qu'il ē possible de tendre vne chorde de luth qui face autant de tours en mesme temps, que celle de mouton, et partant elles feront l'vnisson, Et ce qu'il raporte du verrat, qui par sa mort ou castration fait mourir les cochons dans le ventre de la truye, pō confirmer la hayne desd. chordes a besoing de preuue, particulierem\textsuperscript{t} si le verrat ē chastré, ou tué si \struck{loin} loing de la truye, que son cry ne se puisse oüir. Quant aux chordes d'acier ou de leton, la moindre rouille discontinue leur harmonie
\note{« co\textsuperscript{e} le sieur p. 151 » : lu tel quel ; on attendrait un nom après « sieur ».}

p. 5. au p\textsuperscript{r} corollaire

L'on trouue les surfaces de ces petites chordes comparées a la surface de la barre en disant que les surfaces des cilindres esgaux sont en raison soubsdoublée de leurs hauteurs, n'y comprenant point les bases: par exemple 3 pieds de hauteur de la barre aura 277 fois moins de surface, que le fil au petit mestier tiré d'icelui.

\page{29}
\note{Page de gauche (verso) ; aucun chiffre en tête. Les fins de ligne passent dans le pli.}

p. 9. au commencem\textsuperscript{t}

David au paral. l. 1 ch. 23 auoit estably 4000 chantres pō iouer des instrumens et chanter \ill{} louanges de Dieu, Et Salomon chez Josephe liv. 8 chap. 4 fit 40000 instrumens ps. 88 b\ill{} populus qui scit iubilationem. les trois sortes d'instrumens sont au dernier psalme, laudate eum in tympano, chordis et organo, in cymbalis iubilationis

p. 18 Faber loue grandement Joannes Aurelius Augurelles dans sa Chrysopée. mais s'estoit imaginé que l'agrem\textsuperscript{t} des consonances, et le mouuem\textsuperscript{t} d'vne chorde qui tre\ill{} au toucher d'vne autre a l'vnisson ou a l'octaue, venoit de l'esprit vniuersel si magnifi\ill{} parmy les chymistes, au lieu que tout cela ne vient que de l'air

Je m'estonne qu'a la 24\textsuperscript{me} page il se glorifie de conuaincre que la veüe se fait par emission no\add{n} par reception, et que la voix soit ouye plus viste qu'vn au\textsuperscript{re} son, a cause d'vne petit\ill{} de l'ame qui porte la voix; ce qui monstre qu'il n'estoit pas philosophe, car sa force \ill{} vient de la seule volonté de Dieu, aussy bien que l'eau de ialouzie du 5\textsuperscript{me} cha\ill{} des nombres
\note{« a la 24\textsuperscript{me} page » renvoie sans doute au livre dont il vient d'être question, non à l'Harmonie.}

\marginal{Systeme de 24 quarts de tons ou de 25 chordes, ou tout est egal.}
\[
\begin{array}{r|l|l}
1 & 100000 & 200 \\
2 & 97155 & 194\frac{3}{10} \\
3 & 94389 & 188\frac{3}{4} \\
4 & 91706 & 183\frac{1}{21} \\
5 & 89090 & 178\frac{1}{6} \\
6 & 86559 & 173\frac{3}{25} \\
7 & 84090 & 168\frac{1}{6} \\
8 & 81698 & 163\frac{3}{8} \\
9 & 79369 & 158\frac{3}{4} \\
10 & 77107 & 154\frac{3}{16} \\
11 & 74915 & 149\frac{23}{25} \\
12 & 72782 & 145\frac{3}{8} \\
13 & 70711 & 143\frac{3}{4} \\
14 & 68709 & 137\frac{3}{8} \\
15 & 66742 & 133\frac{12}{25} \\
16 & 64834 & 129\frac{1}{2} \\
17 & 62996 & 125\frac{124}{125} \\
18 & 61199 & 122\frac{12}{31} \\
19 & 59461 & 118\frac{3}{4} \\
20 & 57762 & 115\frac{16}{31} \\
21 & 56126 & 112 \\
22 & 54513 & 109\frac{14}{42} \\
23 & 52973 & 105\frac{115}{125} \\
24 & 51466 & 102\frac{22}{31} \\
25 & 50000 & 100
\end{array}
\]
\note{Le titre et la table sont dans la marge de gauche, sur toute la hauteur des remarques « p. 33 » à « p. 45 » ; le titre est encadré. Les fractions sont petites et quelques-unes restent douteuses (lignes 4, 13, 22, 23).}

p. 33 au commencem\textsuperscript{t}.

Le rencontre de ces deux chordes peuuent peut estre seruir pō trouuer pourquoy, vne mesme chorde fait son propre ton et puis la douziesme, car estant diuisée en 4, la totale fa\ill{} auec sa moitié l'octaue, et auec le quart la 15\textsuperscript{me}, et 3 parts font contre vne part la douziesme

p. 40 au commencem\textsuperscript{t}.

Il y a ce systeme qui ē a la marge

voyez la page 171 des genres et la 439 de la Rythmique ou ce systeme est mis en no\ill{} si l'on prend vn compas de proportion diuisé suiuant la seconde colomne commenceant \ill{} 200, 194 donnera la p\textsuperscript{re} touche pō le quart de ton, et 188 pō le demy ton, et ainsy des autres. ou si l'on commence en bas par 100, 102 donnera le p\textsuperscript{r} quart de ton, et 105 la touche pour le p\textsuperscript{r} demyton.

p. 45. au commencem\textsuperscript{t}.

l'vnisson, ou tel interualle qu'on voudra, se trouue sans oreille, par le nombre des batte\ill{} par exemple si l'on pouuoit tellem\textsuperscript{t}. bander vne chorde de batteau ou vn cable qu'il battist l'air, ou qu'il tremblast cent fois dans vne seconde il feroit l'vnisson a\ill{} chorde de luth qui tremble cent fois en mesme temps, et s'il n'en faisoit que 25 il feroit la quinziesme en bas auec lad. chorde Et cæt.

\marginal{\ill{}iure second des Instruments}
\note{Titre courant dans la marge de gauche, sous la table, relié par une accolade à la remarque qui suit ; la première lettre est dans le pli.}

p. 45. prop. 1\textsuperscript{re}

Les Italiens nomment la seconde figure de la page suiuante Arciluto, on peu\add{t} l'appeller luth a double manche, parce que le tuorbe ē plus grand, et n'a qu\ill{} corde a chasque rang, et n'y a que 30 ou 40 ans que le Bardella l'inuent\ill{} Florence voyez son accord pag. 88. / la figure du luth ē mal tournée

p. 46. au commencem\textsuperscript{t}.

Il y a moyen de f\textsuperscript{e} vn luth a deux corps, auec lequel on iouera tout a vuide auec les deux m\ill{} ce qui a desia esté pratiqué

L'almerie du Maire qui a 15 rangs de chordes, et deux cheualetz, l'vn pō les 4 derniers ra\ill{} \struck{de chordes}, et l'autre pō les 11 derniers se raporte au luth, le manche en est diuisé par quarts de ton afin d'y iouer les 3 genres, ceste diuision se fait en trouuant premierem\textsuperscript{t} 11 moye\ill{} proportionnelles entre les deux lignes qui sont en raison double, et puis en prenant vne m\ill{} proportionnelle entre chascune des 11 precedentes. Il y a 3 sautereaux sous le manche qu\ill{} doibvent toucher par le moyen du pouce de la main gauche 33 touches

voyez le 6. article de la preface generale pō mettre les touches au manche de luth.

p. 47. au commencem\textsuperscript{t}.

L'on peut iouer a dix rangs tout a vuyde, en accordant le luth suiuant ces interualles 3. 4. 5. 6. 7. 8 \ill{} de sorte que de la chanterelle a la 2\textsuperscript{de} il y ayt l'interualle de 12 a 11, et ainsy des autres iusq\ill{} la dixiesme, ou derniere mais il faut que les 4 ou 5 cordes plus aigues soyent quasi toutes d\ill{} mesme grosseur, et partant les chordes marquées par lettres seruiront de tablature voicy \ill{} chanson pō exemple, supposant que les chordes se marquent par a, b, c et cæt. iusques a \ill{} pō la chanterelle

\begin{quote}
hgddfefdeefed

dkgghffdedfgh

hgkdgfghikihg

gdkghfdefghik
\end{quote}
\note{La chanson est écrite en tablature dans la marge de gauche, en quatre lignes de lettres surmontées de signes de valeur (traits, points, rondes) ; les lettres sont données seules.}

\[
\begin{array}{c|c|c|c|c|c|c}
4\cdot 1 & 3\cdot 1 & 1\cdot 2 & 2\cdot 3 & 3\cdot 4 & 4\cdot 5 & 5\cdot 6 \\ \hline
a\,K & a\,g & a\,d & b\,d & a\,b & b\,c & c\,d \\
 & b\,K & b\,f & e\,K & d\,f & f\,g & h\,K \\
 & & c\,h & & g\,K & & \\
 & & d\,K & & & &
\end{array}
\]
\note{Tableau encadré intitulé « consonantiæ », sous le texte. La lettre c a ici la forme d'un r.}

Un autre accord pour neuf chordes 9. 10. 11. 12. 13. 14. 15. 16. et si l'on y ioint 8. 7. 6. 5. \ill{} l'on aura celuy de 13 chordes qui comprend le disdiap\ill{} et que l'on pretend estre celuy de pythagore, et de\add{s} anciens. vne \uncertain{archimole} de 13 chordes est \ill{} pō experimenter ceste sorte d'harmonie
\note{Ce dernier alinéa est écrit à droite du tableau « consonantiæ », séparé de lui par un double trait.}

\page{30}
\note{Page de droite (recto), chiffrée « 15 » à l'encre en haut à droite.}

p. 54. lig. 16. depuis la fin

Ayant pris la moitye de la longueur du luth, depuis le cheualet iusques au sillet, AB, et diuisé en huit partyes esgales, la perpendiculaire DE tirée au milieu, il faut descrire le quart de cercle AC, et du centre B la 6\textsuperscript{me} partye du cercle, AD, laquelle estant diuisée en 12 partyes esgales, si l'on tire des 11 points des lignes au centre du quart de cercle E, ce quart \struck{de} sera tellem\textsuperscript{t}. diuisé en douze partyes inegales que les chordes ou soustendantes entre chasque deux points dud. quart donneront la grandeur des douze touches du manche du luth, dont la 1\textsuperscript{re} plus proche du sillet sera A$\beta$, la 2\textsuperscript{e} $\beta\gamma$ et ainsy des autres. Ceste inuention est de m\textsuperscript{r}. de \uncertain{Forte maison}, neantmoins ces 12 lignes n'estans pas continuellem\textsuperscript{t}. proportionnelles, ceste diuision ne peut estre iuste.
\note{Dans la marge de gauche, la figure : un angle droit en E, la perpendiculaire ED montant jusqu'à D, la base EB graduée ; un arc de cercle part de D vers le bas à gauche et porte deux séries de petites divisions, la seconde près de la lettre C. La lettre A n'est pas visible sur la figure.}

p. 68. sur la fin

voyez vne autre maniere de trouuer les 2 moyennes proportionnelles geometriquem\textsuperscript{t} au dernier liure de la geometrie de m\textsuperscript{r} des Cartes par le moyen d'vne seule parabole et du cercle, ou il demonstre aussy la trisection de l'angle par mesme moyen.

\marginal{\ill{}iure 3\textsuperscript{me} des \ill{}truments a chorde}
\note{Titre courant dans la marge de gauche ; les premières lettres des deux lignes sont dans le pli.}

p. 101. au commencem\textsuperscript{t}

Pline veut qu'on abate le bois au defaut de la lune, lors que le soleil est a la fin du \ill{} ou au commencement du \ill{}, et qu'il sera plus dur coupé lors qu'il gele, et que la lune est sous l'orizon, ce qui rendra le bois de longue durée, et en bon estat. mais les artisans disent qu'il le faut coupper en may, en pleine lune, et qu'il faut mettre le sap et l'erable en quartiers, et le faire seicher au four, 5 ou 6 fois et puis le mettre dans la caue ou lieu moite, afin qu'il prenne vne permanence immobile: ou bien on le met dans la cheminée loing du feu, ou il acquiert vne parfaite coction, qui le rend dur et persistant en mesme assiete, et le despouille de viscosité et verdeur, afin qu'il ne se renfle ny retire point. Faber p. 159 et 169 que le fer s'amolit trempé dans l'eau de fleurs de feue, et s'endurcit dans le ius de l'herbe volubilis: que les graueures sur cuyure se font auec vinaigre \ill{} et sublimé, ou arsenic, ou auec vitriol sel et vinaigre. J'ay veu grauer a Mons auec le seul sublimé sur vn cousteau.
\note{Les deux \ill{} de la première phrase sont deux signes du zodiaque, non des mots : le premier ressemble à une croix de saint André surmontée d'un trait (le Sagittaire ?), le second à un z bouclé (le Capricorne ?).}

\[
\begin{array}{ll}
6^{\text{tt}} & 6\cdot 0\cdot \frac{1}{256} \\
5^{\text{tt}} & 1\cdot 0 \\
4^{\text{tt}} & \\
3^{\text{tt}} & 2\cdot 0\cdot \frac{1}{28} \\
2^{\text{tt}} & 13\cdot 0\cdot \frac{9}{256} \\
2^{\text{tt}} & 4\cdot 0 \\
1^{\text{tt}} & 12\cdot 0\cdot \frac{1}{36} \\
1^{\text{tt}} & 9\cdot 0\cdot \frac{5}{128} \\
1^{\text{tt}} & 4\cdot 0\cdot \frac{1}{64} \\
1^{\text{tt}} &
\end{array}
\]
\note{Table de poids dans la marge de gauche, sous un signe de renvoi en forme de dièse, que l'on retrouve devant la rubrique « p. 123 » qui suit. Le signe « tt » en exposant après le premier nombre est celui de la livre ; les deux nombres suivants sont sans unité écrite.}

p. 123 au commencem\textsuperscript{t}

Faber page 179 calcule aussy les poids pō bander les cordes d'esgale longueur et grosseur, pō f\textsuperscript{e} les degrez diatonics

voyez nos nombres pō les poids qui font tous les interualles harmoniques sur vne mesme chorde liure 3\textsuperscript{me} des mouuements p. 185.

p. 151. prop. 19.

Galilée pag. 103 de ses dialogues se trompe, car la cause pourquoy le son des chordes d'or est plus bas, c'est parce qu'il est plus mol, et non parce qu'il est plus pesant, et la pesanteur d'vn corps ne resiste pas dauantage a la vistesse de son mouuem\textsuperscript{t} que la grosseur.

p. 152. ligne 13 a fine.

Galilée remarque dans son 1\textsuperscript{er} dialogue que la corde d'or esgalle a celle de leton fait la quinte en bas auec elle et qu'vn clauecin tendu de ces deux sortes de chordes sera a la quinte, je trouue plustost qu'ilz font le triton, co\textsuperscript{e} l'on voit en ces experiences marquées par notes, lequel diuise la raison double en deux raisons esgales.

Il faut essayer si vne chorde de leton recuit tendüe de mesme force et estant de mesme grosseur et longueur que la crüe fera mesme son, Et combien plus, et si elles auront mesme raison entre leurs sons, que deux cylindres de mesme matiere egaux, l'vn crud et l'autre recuit

p. 170. lig. 23.

neantmoins lors qu'on aplique sous la teste de ces harpions de petits morceaux de velin, laine ou coton, ilz font vn petit bruit agreable, co\textsuperscript{e} le doux tremblant d'vne orgue, ce que quelques vns prisent beaucoup

p. 176. au commencem\textsuperscript{t}

Faber page 191 appelle la musique de ces bois neutrale et extraordinaire, et dit qu'estant de mesme grosseur ou quarreure, mais differents en longueur suiuant les raisons des interualles, leur quantité et par consequent pesanteur, est comme les nombres qui suiuent $*$ Car il veut que les bastons, qui sonnent vt, RE, soyent co\textsuperscript{e} 18 et 17 racines imaginables de 8 et 9, parce que 17 multiplié par soy c'est a dire 288, fait auec 18 multiplié par soy c'est a dire 324 vne raison sesquioctaue et ainsy des autres nombres de la 2\textsuperscript{e} colomne qui sont \struck{qui se} les quarrez des nombres de la 1\textsuperscript{re} colomne, la 3\textsuperscript{me} colomne contient les nombres lineaires des interualles harmoniques commenceant par 18 et la 4\textsuperscript{me} contient leurs quarrez.
\[
\begin{array}{|l|l|l|l|}
\hline
12\frac{3}{4} & 162 & 9 & 4\frac{1}{2} \\
13\frac{1}{2} & 182\frac{1}{4} & 10\frac{1}{8} & 5\frac{89}{128} \\
13\frac{23}{27} & 192 & 10\frac{2}{3} & 6\frac{26}{81} \\
14\frac{5}{7} & 216 & 12 & 8 \\
15\frac{3}{5} & 243 & 13\frac{1}{2} & 10\frac{1}{8} \\
16\frac{1}{8} & 256 & 14\frac{2}{9} & 11\frac{19}{81} \\
17 & 288 & 16 & 14\frac{2}{9} \\
18 & 324 & 18 & 18 \\ \hline
\end{array}
\]
\note{Table dans la marge de gauche, sous un astérisque qui répond à celui du texte (« qui suivent $*$ »). « 17 multiplié par soy c'est a dire 288 » : lu tel quel ; $17^2 = 289$, et 288 est le carré de la « racine imaginable » (irrationnelle) que 17 représente.}

Il faut examiner si ces bois de mesme grosseur et de diuerse longueur font les tons proportionaux aux tuyaux d'orgue de mesme grosseur, et differens en longueur.

\page{31}
\note{Page de gauche (verso) ; aucun chiffre en tête. Les fins de ligne passent dans le pli. Dans la marge de gauche paraît le bord de deux figures qui sont sur un feuillet dépliant, relié contre cette page et photographié déplié à la vue 32 ; elles sont décrites à cette vue.}

p. 211. coroll. 2\textsuperscript{me}

Il faut voir dans le 1\textsuperscript{er} dialogue de Galilée, ou dans mon 21 article de sa version, ce qu'il \ill{} des ondes fremissantes d'vn verre, qu'il dit se fendre en deux, quand le son monte a l'octa\add{ve} voyez la prop. 17 du 7\textsuperscript{me} liure suiuant p. 32 et 33.

p. 218 au commencem\textsuperscript{t}

Quand sur la chorde MN, le doigt est mis sur C la totale MN resonne contre MC, la quart\ill{} contre CG la quinziesme, partant la 15\textsuperscript{me} peut resonner quand on entend la quarte et vice versa, ainsy le C et G resonneront: et parce que D graue de MN, se trouue auec d aigue de PQ, la chorde peut resonner la quinte, la douziesme et la quarte co\textsuperscript{e} complem\textsuperscript{t}. de l'octaue: et c'est au\ill{} complem\textsuperscript{ts} de consonances qu'il faut prendre garde. a et F se trouuent aussy ensemble et p\ill{} sonner la 9\textsuperscript{me} ou la 7\textsuperscript{me} voyez Faber pag. 123.
\note{Les lettres M, N, P, Q, C, D, G renvoient à la figure du monochorde, deux échelles graduées sur le feuillet dépliant de la vue 32.}

Il faut examiner si ces correspondances de chordes peuuent estre cause des 4 ou 5 sons que fait \ill{} mesme chorde: et si cela ne peut satisfe\add{re}, il y a aparence que les mouuem\textsuperscript{ts} des sons aigus resonans, viennent des mouuem\textsuperscript{ts} de certaines partyes internes de la chorde.

Le cheualet ē \uncertain{q\textsuperscript{q}} apellé par Faber aphonius, a raison qu'il est co\textsuperscript{e} l'interualle de l'infinité \ill{} sons, qui sont en la chorde NM

Pour trouuer toutes les diuisions de son monochorde il diuise 1\textsuperscript{o} la chorde GM en 4 p\textsuperscript{ties} e\ill{} qui donnent GC, et G et g / 2\textsuperscript{o} CM en 4 p\textsuperscript{ties}, d'ou vient F, c et c / 3\textsuperscript{o} FM en quatre par\ill{} d'ou viennent pō auoir la feinte dA, et F, et f puis il diuise aXM en \ill{} qui donnent le\add{s} feintes entre D et E / 4\textsuperscript{o} DXM en 3 pō auoir les feintes de G / 5\textsuperscript{o} il diuise G ou NM en 3 pō auoir les deux D et d, et puis diuise DM en 3, pō auoir A, et puis AM diuisée \add{en} 3, donne E, et consequem\textsuperscript{t}. il faut diuiser EM en 3 pō auoir B, et finalem\textsuperscript{t} partir BM en 3 pō av\add{oir} la feinte de F. Quelques practiciens marquent les diuisions des chordes ainsy 442 \ill{} ou 33 R. 33 et cæt. pō dire qu'il faut premierem\textsuperscript{t} diuiser en 4 p\textsuperscript{ties}, puis encore en 4 et puis en 2, puis retrograder de l'autre costé voyez Faber pag. 132
\note{« divise aXM en » : un blanc est laissé après « en », avant « qui donnent ». La lettre suivie d'une croix (aX, dX, fX…) note la feinte, c'est-à-dire le dièse de la note. « R. » est l'abréviation que le copiste emploie ailleurs pour « racine » ; ici elle semble désigner un renvoi d'une division à l'autre.}

4 signifie GCgg / 4 depuis C marquent f cc / 2 f en haut / R. de lad. F donne aX /

4 d A'X feinte donne dX aX, aX, / 2 font dX dX / R. fait gX / 4 produit cX gX gX / 4 produi\add{t} fX, cX, cX. / 2 fait octaues fX, fX / R. fait B / 4 produit eX, BB / 2\textsuperscript{o} pour les doubles / R. \ill{}

4. d. a a. / 2 d. d. / R. g. voyez Faber page 137. de mesmes par la diuision en 3.

3 fait G dd / d en 3 fait aa / a porté par R. donne E / E mis en 3 BB / B en 3 FX fX / fX en 3 cX c\ill{} c en 3 gX gX / gX en \uncertain{8} dX dX / lequel retrogradé fait aX aX / aX fait f f / et f R. C / et C en 3 g g / et g en 3 fait d. d. restent les repetitions pō les octaues /

Il nomme Cæsar Cæsarianus commentateur de Vitruve qui traicte du Chromatiq au 5. l\add{ivre} de son Architecture chap. 5 et \uncertain{sert} seulem\textsuperscript{t}. quasi de Gaffurus et Frisichius qui escriuit l'an \ill{} et que Gafurus fut chantre de l'Eglise de Milan vers l'an 1484 et musicien de Fran\ill{} premier, et que Louys Sforce en estoit duc. /

p. 224. lig. 8 a fine

voyez ceste preuue mieux expliquée p. 51, 52 et cæt. du liure de la balistique ou la mesme ligne ē coupée 12 fois en moyenne et extreme raison.

p. 225 lig. 5.

Il est vray que le moindre segment osté du plus grand, a scauoir BE de EA; EF demeure le mo\ill{} et led. moindre BE, deuient le plus grand, et ainsy des autres dont j'ay parlé dans la prop\ill{} de la balistique, et partant Salinas ne manque pas en cela, mais en ce qui est dit icy d\ill{} qu'vne ligne ainsy coupée en 12 partyes continuellem\textsuperscript{t} proportionnelles donne les 12 demy\ill{} ce qui ē faux.

p. 228. au commencem\textsuperscript{t}.

L'on pourroit icy adiouster la canne harmonique qui ē venüe de Cambaye Royaume de \ill{} entre le Sein persique et Arabique, dont les cordes sont faites de la mesme canne, a la\ill{} elles sont continuées, auec de petits cheualetz \ill{}, et de petites mortaises dessous, d'ou elles ont esté tirées.

Ces grosses courgues instrum\textsuperscript{ts}. de la chine sont de boys fort mince pō augmenter l'harmonie de la canne, AE, a laquelle on peut adiouster vn tampon et vne lumiere pō seruir de fleute: et si l'ouuerture des courgues estoit couuerte de parchemin, elles pourroient seruir de tambour. Voyez d'autres instruments des Indes, en l'histoire orientale Et occidentale, et voyez aussy le liure de Pratorius des instruments.

p. 228. au commencem\textsuperscript{t}.

Faber page 31 veut que les flustes ayent leur longueur en mesme raison de l'interualle de musi\add{que} qu'elles font, leur largeur soit en raison des racines des termes dud. interualle, par exem\add{ple}
\note{Au pied de la page, à droite, la réclame « si elles » ; la phrase se poursuit à la vue 33, la vue 32 étant le feuillet dépliant.}

\page{32}
\note{Vue d'un feuillet dépliant, de format plus grand que les pages, relié contre la vue 31 et photographié déplié ; le reste de la vue est la face blanche d'une page. Ce sont les deux figures dont on voit le bord dans la marge de gauche de la vue 31. En haut à droite, une bande verticale porte le monochorde : deux échelles graduées, lettrées de M à N et de P à Q (voir la note de la vue 31, remarque « p. 218 »). En bas, une feuille collée ou dépliée porte quatre instruments à vent dessinés debout, chacun gradué de 1 en bas à 27 en haut.}

Liure 5\textsuperscript{me} des instruments a vent propo\textsuperscript{on} 6\textsuperscript{me} p. 232
\note{Titre en tête de la feuille des instruments à vent. Les quatre figures, de gauche à droite, légendées verticalement : « Cornet a bouquin », tuyau courbe à embouchure en cuvette, gradué de 1 à 27 ; « phifre ou fluste traversine », tuyau droit, « 27 embucheure » en tête, des trous marqués d'un cercle pointé aux degrés impairs de la moitié inférieure ; « Fluste douce », tuyau droit à pied évasé, « 27 bouche » en tête, des trous marqués d'un cercle pointé, celui du degré 18 en pointillé ; « Flageolet », tuyau court à pied évasé, les degrés 27 et 28 dans un petit carré en tête, et en regard des trous 5, 7, 9, 12, 16 et 17 les mesures $4\frac{1}{2}$, $6\frac{3}{4}$, 9, $12\frac{1}{3}$, $15\frac{3}{4}$, $16\frac{3}{4}$.}

\page{33}
\note{Page de droite (recto), chiffrée « 16 » à l'encre en haut à droite.}

si elles sont en longueur de 9 a 8, il veut que \struck{le} \add{leur} contenu ou surface soit co\textsuperscript{e} la R. de 9 a 8 mais parce que 8 n'a point de racine il prent la moitie de la raison sesquioctaue a peu pres, a scauoir de 17 a 18, dont les quarrez sont 324 et \struck{\ill{}} 289. de sorte qu'il veut que 17 et 18. soyent R. de sesquioctaue moins $\frac{1}{288}$, de sorte qu'il faut souffrir l'excez de $\frac{1}{32}$ de sorte que la flute vt sera co\textsuperscript{e} 81 en son contenu et Re co\textsuperscript{e} 64.
\note{Suite de la remarque « p. 228 » de la vue 31 (réclame « si elles »). « leur » est écrit au-dessus de « le », qui n'est pas nettement biffé.}

Il luy eut esté bien plus aisé de proceder par moyennes proportionelles ou par les 3 proportionelles, ou bien mesme de f\textsuperscript{e} que la largeur des flutes, c'est a dire le diametre de leurs bases, eussent esté en mesme raison que leur longueur

Puisque la circonference de 22 donne $38\frac{1}{2}$ en son aire, et le quarré isoperimetre n'a que $30\frac{1}{4}$, il faut voir si deux fleutes de mesme longueur l'vne ronde et l'autre quarrée isoperimetres qui sont entr'elles co\textsuperscript{e} 154 a 121, qui ē quasi vne sesquiquarte, co\textsuperscript{e} 14 a 11 different de son et combien

On peut aussy f\textsuperscript{e} vn tuyau triangulaire isoperimetre des precedents qui ayt $7\frac{1}{3}$ de chasque costé de son aire de $23\frac{41}{135}$ ou quasi $\frac{8}{27}$ de sorte que sa capacité sera a celle du cercle co\textsuperscript{e} 77 a 46, et voir sa difference de son auec les deux precedentes.

ce que l'on peut resoudre par ce qui ē dit dans la propo\textsuperscript{on} \quad du liure des orgues
\note{Un blanc est laissé pour le numéro de la proposition.}

p. 232. propo\textsuperscript{on} 6\textsuperscript{me}

voyez le dessein d'vn flageolet, d'vne flute douce, d'vn phifre et d'vn cornet a bouquin cy attaché.
\note{Ce sont les figures du feuillet dépliant photographié à la vue 32.}

p. 250. au commencem\textsuperscript{t}

Ce phenomene des instruments a vent ē difficile a expliquer; car pō peu de vent qu'on leur donne ilz passent a l'octaue, et semblent se batre sur ce point. L'on peut donc dire que quand le tuyau, AB, reçoit le vent naturel il altere son air entier et fait son 1\textsuperscript{er} son graue, de sorte que l'air qui est en B, forme le son, car si celuy qui ē en, C, auoit iour, il donneroit vn autre son
\note{Dans la marge de gauche, la figure : un petit tuyau vertical, A en haut, C au milieu, B en bas.}

Et puisque toute la force du vent repond a toute la force qui meut le poids en la balance, dans laquelle si l'on aiouste, ou si l'on diminüe, de la force, quand ilz sont egaux, l'vn ou l'autre l'emporte, de mesme qu'aux sons lors qu'ilz passent de l'vn a l'autre. De mesme si l'on veut encore adiouster vn son plus aigu au graue precedent, il faut alterer tout l'air qui est depuis, A, iusques a, B, puis que le son se fait en, B, qui est de mesme qu'vn autre poids esgal adiousté au p\textsuperscript{er}: D'ou il arriue que tous les sons de la trompette suiuent l'ordre naturel des nombres 1, 2, 3, 4 et cæt. parce que la nature ne peut faire de sons dans les instrum\textsuperscript{ts}, qu'en alterant tout leur air autant de fois que les sons se changent. Et cela se trouue par experience iusqu'au nombre, 6, mais le reste ne se faisant pas ainsy, puisque la trompette saute de 6 a 8, il faut examiner si ses sons ne se font point par l'industrie de l'embouchure, autrem\textsuperscript{t} ceste raison ne contente non plus que celle qui est expliquée en ceste proposition, ou celle qui est dans l'harmonie latine liure 2\textsuperscript{on} des instruments propo\textsuperscript{on} 18\textsuperscript{me}. mais ceste mesme raison est descritte a la fin des obseruations

p. 283 au commencement.

Il faut remarquer que les tuyaux ou chalumeaux de la chalemie, cornemuse et musette, n'octauient pas co\textsuperscript{e} le flageolet et la flute, quoy que certains instrum\textsuperscript{ts} plus grands ayent des hanches, et que l'on ne laisse pas certains trous ouuerts en jouant co\textsuperscript{e} on fait au flageolet de sorte qu'il est besoing d'expliquer co\textsuperscript{e} l'on bouche, et co\textsuperscript{e} l'on ouure les trous pour faire chasque son

\marginal{\ill{} \uncertain{vij}\textsuperscript{me} liure instruments de \ill{}vssion}
\note{Titre courant dans la marge de gauche ; le début est dans le pli. Sans doute « Livre … des instruments de percussion », que la vue 27 cite (« livre des instrumens de percussion »).}

p. 6. au commencem\textsuperscript{t} et 7\textsuperscript{e} ensuiuant

L'on void encores vne autre chose icy attachée laquelle estant bien proportionée pour seruir d'Idée et de modelle aux fondeurs. or elle a d'vne pince a l'autre de, B, a, A, 15 fois l'espesseur de son bord. sa hauteur iusques au dehors du cerueau ē de $12\frac{1}{2}$ \struck{et} son diametre ou commence la courbeure du cerueau 8: la haute perpendiculaire iusqu'au dit diametre $11\frac{1}{2}$ sa montée 12: le cerueau superieur a 5 bords en son diametre: Du fort du bord iusqu'au bord de la pince tant exterieurem\textsuperscript{t} qu'interieurem\textsuperscript{t}, $1\frac{1}{2}$. l'espesseur aud. fort du bord, 1, a la moitie ou faussure de la cloche $\frac{1}{3}$ de bord, et par tout autant en montant depuis la faussure comme aussy au cerueau superieur
\note{La « chose icy attachée » est le dessin d'une cloche ; il n'est pas dans cette vue (voir la vue 34).}

\page{34}
\note{Un feuillet volant de format réduit, portant le dessin d'une cloche, photographié posé sur la page suivante (dont on voit le bord droit, écrit, à droite de la vue ; ce texte est lu à la vue 35). C'est la « chose icy attachée » de la dernière remarque de la vue 33. Aucun texte, aucun chiffre de page sur le feuillet.}

\note{Le dessin, à la plume, montre la cloche en coupe, les parois hachurées, surmontée de ses anses (le cerveau et la couronne, lettrés G et F), avec la construction en pointillés qui sert à en tracer le profil. Lettres lues : A et B aux deux extrémités de la base (ligne A I Q o p Z T u B) ; aux bords, g, K, M, L, e ; sur l'axe vertical, de haut en bas, Y (?), \&, f, h, s, O, N, I ; au niveau de la faussure, G i … f … q X H ; en haut, D m … Y E, d, n, n, c, c ; à droite le point R, d'où partent des lignes vers P et vers le bord ; au milieu, P. Les lettres de la construction sont celles que nomme le texte de la vue 35.}

\page{35}
\note{Page de gauche (verso) ; aucun chiffre en tête. Les fins de ligne passent dans le pli et sont souvent perdues. À gauche de la vue paraît le bord du feuillet de la cloche (vue 34). Une grosse tache d'encre ronde couvre un mot vers le milieu de la page.}

que les anses doibvent quasi occuper. la maistresse anse doibt auoir $2\frac{1}{2}$ de hauteur, et de grosseur $\frac{2}{3}$
\note{Suite de la remarque « p. 6 » de la vue 33, sur les proportions de la cloche.}

\subsection*{Construction}

Soit menée la ligne AB de 150, supposant chasque bord diuisé en 10 parties: et, AB, estan\add{t} diuisée en 2 parties esgales au point I de chascune 75. dud. point soit esleuée la perpendiculaire, IC, sur AB, puis prenant 125 sur IC, a commencer du point I, vers C, \ill{} fait le point, F, qui sera la hauteur perpendiculaire de la cloche. puis soyent prises 115 partyes du point I, vers C, et soit fait le point \&. par lequel soit menée la ligne DE, parallele a, AB, puis de part et d'autre dud. point \&, soient posées 40 partyes \ill{} terminants aux points, D et E, et soyent menées les lignes BE, et AD, chascune de 120. \ill{} soit diuisée la ligne, \&I, en deux esgalem\textsuperscript{t} au point f, par lequel soit menée la ligne GH \ill{} parallele a AB, puis soit descrit l'arc IR de l'interualle, BI, coupant la ligne, BE, au point, R, duquel soit menée la ligne RI coupant la ligne GH. en X, puis soyent portées 15 partyes du point, B, sur la ligne BE, et soit fait le point K, duquel co\textsuperscript{e} centre et de l'interualle BK soit descrit l'arc KZ. puis soyent portées 50 partyes du point R sur l'arc, RI, et soit fai\add{t} le point P, par lequel soit menée la ligne BPh, coupant l'arc KZ en L, et soit menée la ligne KL, ce que l'on peut repeter sur le costé DA, mais il suffit du point h, duquel l'on peut mener la ligne Ah, puis les lignes KK, et LL, paralleles a AB, pō auoir les sections KML du costé A semblables au costé B. soit aussy portée l'ouuerture fX, de f vers G, et soit fai\add{t} le point i puis soit prise l'ouuerture DE, ou 80 partyes, et icelle portée de F vers I, et soit fait le point s, duquel co\textsuperscript{e} centre, et de la mesme ouuerture soit descrit l'arc DFE, passant precisem\textsuperscript{t} par F. apres soit diuisé KL en trois p\textsuperscript{ties} esgales dont l'vne soit portée du point X en H, et soit fait led. point, H, co\textsuperscript{e} aussy semblablem\textsuperscript{t} de i en G, de E en Y et de \ill{} V et de D en m. puis par la 25\textsuperscript{e} du 3\textsuperscript{e} soit cherché le centre d'vn cercle passant par les \ill{} points LXY, et par miL, desquels points soyent aussy descrits les arcs DG et HE, ou autrem\textsuperscript{t}. prenant 4 fois la ligne PR ou 200 parties, soyent du point Y, et L, descrits deux arcs a main droitte \ill{} hors de la figure, et ou se fera l'intersection d'iceux ce sera le cen\add{tre} duquel l'on pourra \ill{} descrire lesd. arcs LXY, et EHK: mais prenant cent partyes et des points H, K descriuant co\textsuperscript{e} dessus deux arcs, le point de l'intersection sera le centre, duquel co\textsuperscript{e} deuant \struck{descrire} l'on pourra descrire les arcs HK et GK. et finalem\textsuperscript{t} descriuant du poin\add{t} s, l'arc mVY, le corps de la cloche sera tracé, excepté le cerueau superieur, qu'il faut \ill{} de la distance du point F, vers I, de six partyes du diametre susd. et pō la corbeure elle \ill{} descrit de l'interualle de 15 partyes, depuis le point, E, ou D, les points n. seruants de cen\add{tre}
\note{La « 25\textsuperscript{e} du 3\textsuperscript{e} » est une proposition du troisième livre d'Euclide ; le copiste ne nomme pas le livre. Le « \&I », le « point \& » : le signe \& sert de lettre à la figure (vue 34, sur l'axe). « s » et « f » sont les lettres minuscules de la figure ; certaines, lues sur la seule écriture, restent douteuses (« miL »).}

Mesure de chascune des lignes de ceste figure

\[
\begin{array}{ll}
\text{la ligne DE} & 80 \\
mY & 73\frac{1}{3} \\
GH & 41\frac{5}{6} \\
iX & 85\frac{7}{40} \\
YX \text{ ou } \& f & 57\frac{1}{2} \\
KOK & 141\frac{5}{9} \\
MOM & 117\frac{8}{9} \\
OF & 43\frac{1}{5} \\
LNL & 125\frac{1}{5} \\
ON & 6 \\
AB & 150 \\
NI & 8\frac{9}{10} \\
KM \text{ ou } ONI & 14\frac{3}{10} \\
Bu & 4\frac{3}{10} \\
BT & 12\frac{2}{3} \\
BP & 32\frac{33}{80} \\
Bo & 35 \\
PR & 50
\end{array}
\]

\[
\begin{array}{lr}
\text{Le cerueau contient de solide pō le corps DFEVY} & 15058 \text{ lignes} \\
\text{Ses filetz cc} & 157 \text{ lignes} \\
\text{Les quatre anses} & 8280 \\
\text{Le solide du corps DEGH} & 333443 \\
\text{Son vuyde mYXi} & 283766 \\
\text{Solidité du metal DmGi\&} & 49677 \\
\text{Le solide du corps GHKK} & 462203 \\
\text{Son vuide iXMM} & 348760 \\
\text{Solidité du metal GiKM et XHMK} & 113443 \\
\text{Les filetz d, d.} & 254 \\
\text{Les 3 filetz ee,} & 612 \\
\text{Le solide du corps KKBA} & 238770 \\
\text{Le vuide MMLL} & 68205 \\
\text{L'autre vuyde LLAB} & 124971 \\
\text{Et pō les deux vuydes} & 193176 \\
\text{Solidité du metal KMMK et ALLB} & 45594 \\
\text{Les deux filetz gg} & 475 \\
\text{Fait en tout de solide pō tout le corps de la cloche} & 233550 \text{ lignes}
\end{array}
\]
\note{Les deux tables sont côte à côte sur la page ; dans la seconde, les articles sont reliés à leurs nombres par des lignes de points. « 124971 » : le dernier chiffre est incertain (1 ou c). Les lettres des deux colonnes sont celles de la figure de la vue 34.}

L'aueugle d'Estampes, qui en retirant son haleyne dans vn verre ou cloche leur fait f\textsuperscript{e} les \ill{} tons de la cloche, ou verre, tient que la cloche a l'estendue de 3 octaues, dont la plus basse respond \add{au} bord HI, la 2\textsuperscript{de} au milieu FG, et la plus haute au haut DE, et que la cloche est d'autant meilleure, \ill{} plus mauuaise, qu'elle aproche plus ou moins de ces trois octaues, ou plustost que le plus \struck{\ill{}} \add{bas} est de \ill{} F, iusqu'a D, le tenor de D a B et le dessus de F a H. Et que lors qu'il s'y trouue quelque desfau\add{t} ou mauuais ton cela vient de ce que le metail est trop espais, ou tenu, ou que la cloche ē trop large, \ill{} estroite en quelque lieu: et le son ē vacillant quand la matiere n'est pas bien fondüe, et qu'il y demeu\add{re} des vents, ou choses semblables. les differents traits des diuerses ouuertures du compas qui forment \ill{} moule, peuuent estre la cause de ces sons differens. Voyez la 18\textsuperscript{me} proposition
\note{Dans la marge de gauche, en regard de la ligne qui nomme « FG », les lettres « FG » répétées. Au pied de la page, à droite, la réclame « p. 19 au commencem\textsuperscript{t} », qui est la rubrique de tête de la vue 36.}

\page{36}
\note{Page de droite (recto), chiffrée « 17 » à l'encre en haut à droite. Dans la marge intérieure, contre le pli, quelques mots d'une encre plus noire et d'un trait plus gros, dont le début est pris dans la reliure ; ils sont donnés à leur place, en \textit{marginal}.}

p. 19. au commencem\textsuperscript{t}

Le poids et la mesure d'vne cloche. J'ay pris de l'eau dans vn vaisseau en forme de cylindre de 13 pouces 9 lignes et $\frac{1}{3}$ de diametre: et puis $52^{\text{tt}}\,\frac{4}{5}$ de metal de cloche composé de $\frac{3}{4}$ de rosette fine, et d'vn quart d'estain fin, dans lad. eau, ce metal a fait renfler l'eau de 12 lignes de haut, ce qui fait de solide pō l'occupation dud metail \uncertain{197941} lignes cubes, or 4 pouces cubes dud. metail par obseruation, pesent 13053 grains, qui font 22 onces, et 21 grains. or les 4 pouces valent 6912 lignes, d'ou il s'ensuit qu'il y a 4880 lignes cubes au solide d'vne liure de metal, donc si l'on multiplie 4880 par $52\frac{4}{5}$, il viendra 257664, qui n'est que 123 lignes d'erreur, qui peut auoir esté commis par l'eau. Ayant encores experimenté plus\textsuperscript{rs} cloches de la construction de la cloche mise cy dessus, et de la mesme composition de metal, j'ay trouué que la cloche qui pese $25^{\text{tt}}$ a iustem\textsuperscript{t} 8 lignes et $\frac{1}{25}$ de bord, supposé que le pied sur lequel sont prises lesd. mesures, soit bien rectifié co\textsuperscript{e} le mien: que celle qui pese $48^{\text{tt}}$ a iustem\textsuperscript{t} 10 lignes d'espesseur de bord, et celle qui pese $622^{\text{tt}}$ a 23 lignes $\frac{1}{2}$.
\marginal{\ill{} y a quelque \ill{}vtre erreur en ce cal\ill{}}
\note{« $52^{\text{tt}}$ » : le signe de la livre, deux t liés, est rendu tt en exposant. « 197941 » : lecture douteuse ; les deux derniers chiffres sont serrés et surchargés. Le texte donne ensuite 257664 et « 123 lignes d'erreur », ce qui ferait attendre 257787 ; la note marginale en regard signale une erreur dans ce calcul.}

Prenons la plus reguliere des 3, qui ē de $48^{\text{tt}}$ et 10 lignes d'espesseur de bord, auec le poids susd. des 4 pouces de mesme metail, et pō l'explicaon de la figure cy deuant, nous supposerons des partyes au lieu des lignes, par lesquelles toutes cloches peuuent estre mesurées. Donc la ligne, IB, estant posée de 75, puisque la toute AB, est de 150, et la ligne, \&E, de 40, retranchant, \&E, de, IB, restera, OB, 35, donc nous auons du triangle, EOB, les trois costez cognus, auec l'angle, E,OB, qui est droit par sa construction, partant l'angle, EBO, sera trouué de 75 degr. 3, et l'angle, OEB, de 16 degr. 57, d'ou s'ensuit que du triangle, RBI, les angles BIR et BRI: sont chacun de 53 deg. $28\frac{1}{2}$ faisant donc du point, X, tomber la perpendiculaire XP, sur IB, l'angle IXP sera de 36 deg. $31\frac{1}{2}$ et partant la ligne IP, sera de $42\frac{47}{80}$ de FX son esgale, et partant FH, de $45\frac{11}{12}$. De rechef les 3 costez du triangle isoscele RBD sont cognus, soit donc imaginée vne perpendiculaire tombant de l'angle B, sur la base PR au point q, elle diuisera la base en deux esgalem\textsuperscript{t}, or elle a esté posée de 50, donc la moitye sera 25 et partant l'angle PBR sera de \struck{\ill{}} 38 deg. 36 et l'angle PBI de 34 deg. 7 et chascun des angles BPR, et BRP de 72 deg. 32. le Triangle XVB ē æquiangle au triangle EOB, partant les trois angles sont cognus auec le costé KB, donc le costé XV sera trouué de $14\frac{3}{10}$, et le 3\textsuperscript{me} uB de $4\frac{7}{20}$
\marginal{na ces angles sont faux faux}
\marginal{RBP}
\marginal{\ill{}ois KMB \ill{} y a faute \ill{}ois Ku}
\note{« 75 degr. 3 », « 16 degr. 57 » : degrés et minutes. Les trois notes marginales sont en regard, de haut en bas, des lignes sur les angles BIR et BRI, de « RBD sont cognus », et de « Triangle XVB … KB » ; « RBP » y corrige sans doute « RBD » du texte, souligné. « na » est sans doute l'abréviation de « nota » (même marque à la vue 39) ; « faux » est écrit deux fois, l'un sous l'autre, le premier repassé.}

p. 32. prop. XVII

Il faut experimenter vn verre cylindrique co\textsuperscript{e} CDEF, et marquer les lieux de ses interualles, et aussy les lieux du verre conique CBD, afin de voir combien ilz seront differents de hauteur pō f\textsuperscript{e} les tons de l'octaue

La hauteur du verre s'entend de la hauteur de sa coupe
\note{Dans la marge de gauche, la figure : un rectangle en pointillés CDEF, et dedans un triangle à sommet B en bas, A en haut au-dessus ; lettres lues A, D, B, F.}

p. 33 au commencem\textsuperscript{t}

Donc Kirker se trompe p. 847 de magnetismo musica, quand il dit que le verre demy plein faict l'octaue en bas, contre le mesme verre vuyde. voyez s'il ne se trompe pas encore, d'autant que le verre estant diuisé en 5 partyes esgales et qu'on en remplisse 3, il fera la quinte, au lieu qu'il debvroit descendre plus bas que l'octaue, selon luy, puisque sa moitié pleine qu'il dit faire l'octaue ē moindre que $\frac{3}{5}$.

page 861. il marque que de quatre verres D A C, ainsy disposez, A plein d'eau de vie bouillonne ou fremit plus fort, B plein de bon vin apres, D apres d'eau subtile et finalem\textsuperscript{t} C d'eau plus grossiere fremissant de moins en moins, lors qu'on frote le bord du verre vuyde E, de sorte que que A represente la cholere ignée, B la complexion sanguine, et l'air, D l'eau et le flegme et C, qui a peyne fremit la terre et la melancholie.
\note{Les quatre verres sont figurés dans la ligne par de petits cercles lettrés, disposés en croix : B au-dessus, D, A, C sur la ligne, E au-dessous. « que que » : le mot est répété.}

voyez la remarque de l'aueugle a la marge de la 4\textsuperscript{me} obseruation

ibid. lig. 8 deuant la fin

Ce\uncertain{cy} n'est pas entierem\textsuperscript{t}. conforme a ce qui est en la page precedente vis a vis de cecy, ce qui qui monstre la difference des experiences.
\note{« qui qui » : le mot est répété.}

\page{37}
\note{Page de gauche (verso) ; aucun chiffre en tête. Les fins de ligne passent dans le pli. À droite de la vue paraît le bord d'une figure de la page suivante (vue 38).}

p. 36. lig. 10

car le tonneau demeure ouuert, au lieu que le tuyau d'orgue susd. a demy plein demeure bouché, ce qui empesche la contrarieté que l'on pourroit former icy.

ibid. prop. XVIII

vn aueugle d'Utrecht poussant ou retirant son vent au bord d'vne cloche, ceste cloche resonne a son ton naturel, et a l'octaue, a la douziesme et cæt. il fait la mesme chose auec vn verre. voyez la 5\textsuperscript{me} propo\textsuperscript{on}

Le ton resonant de la cloche est apellé par les fondeurs le Cornet, chez Jean le Fevre secret\textsuperscript{re} du Cardinal de Givry euesque de Langres, en son 3\textsuperscript{me} liure de sa triple musique

p. 37. prop. 19 lig. 18.

l'on peut expliquer ceste dificulté par celle du retour de l'arc, ou du poids attaché a \ill{} chorde, soit par les esprits en los qui courent ça et la, estans accoustumez aux \ill{} ou a la figure de la cloche, soit par le mouuem\textsuperscript{t} de la matiere subtile
\note{« en los » : lu tel quel.}

p. 60. a l'aduertissem\textsuperscript{t}. lig. 17.

l'archet de peau de Cheual auec son poil s'est trouué le meilleur de tous pour \ill{} douceur, mais il ne faut pas qu'il se mousse a contrepoil.

Et l'on met mainten\textsuperscript{t}. deux ieux a l'archimole, qui sont a l'octaue ou a l'vnis\add{son} l'vn de l'autre: si l'on fait des ieux a la quinte et a la tierce, l'on imitera qua\ill{} tout ce qui se fait sur l'orgue.

Et sur le coing de la table l'on met vne chorde auec toutes les diuisions de l'octaue laquelle sert de monochorde pō accorder vne octaue entiere toute a l'vnisson \ill{} 13 touches dud. monocorde touché par l'archet. A quoy l'on peut adiouster \ill{} tambour, ou barillet, qui fera son mouuem\textsuperscript{t} par des ressorts que l'on n'aura \ill{} besoing de toucher; de plus on y adiouste vn tremblant, que l'on en oste quand \ill{} veut, et vn ressort qui tourne la roüe de plomb 2 heures durant, fait co\textsuperscript{e} celu\add{y} d'vne harquebuse, ou d'vne monstre d'horologe.
\note{« archimole » : ainsi écrit, comme à la vue 29.}

\marginal{De l'vtilité de l'harmonie}
\note{Titre courant dans la marge de gauche, sous un trait. Les renvois qui suivent (p. 2, 3, 4…) repartent d'une nouvelle pagination.}

p. 2 lig. 36.

l'experience n\textsuperscript{o} a fait voir que l'esprit de vin auec la vessie ne vaut rien pō ce su\add{jet} car il n'enfle pas et le bout de la vessie estant detaché iette sa graisse auec gran\add{de} incommodité, et la fumée de l'esprit de vin se refroidit incontinent par le contact d\add{e la} vessie quoy que tenüe chaude auec vne fermoire, ou autrem\textsuperscript{t}
\note{« n\textsuperscript{o} » : l'abréviation de « nous ».}

p. 3 lig. 22.

Il est mainten\textsuperscript{t}. aysé de peser l'air dans le vuyde par le moyen du tuyau ou bouteille ou le mercure descendant fait le vuyde d'air.

lig. 29. l'essay fait suiuant ceste methode n'a pas reussy, la vapeur se tourna en eau auant qu'elle enfle la bouteille.

p. 4. lig. 10.

voyez la maniere dont Galilée pese l'air dans le vuyde dans son 1\textsuperscript{er} dialogue ou d\add{ans} le 16\textsuperscript{me} article du 1\textsuperscript{er} liure de n\textsuperscript{re} version ou il trouue qu'il est 400 fois plus leg\add{er} que l'eau.

M\textsuperscript{r}. des Cartes a trouué que l'air est a l'eau co\textsuperscript{e} 1 a 145, auec vne phiole souflée a la la\ill{} de la grosseur d'vne balle de tripot, qui froide pesoit $78\frac{1}{2}$ grains et chaufée 78 grai\add{ns} et refroidie l'eau succée par le bec de la phiole, l'eau pesoit $72\frac{1}{2}$, de sorte que l'air so\add{rti} estoit a l'eau entrée co\textsuperscript{e} $\frac{1}{2}$ a $72\frac{1}{2}$ qui est co\textsuperscript{e} 1 a 145.
\note{« refroidie » : « par » est écrit au-dessus de ce mot. « souflée a la la\ill{} » : sans doute « a la lampe », la fin du mot étant dans le pli.}

voyez la 29. prop. de la pneumatique, ou l'air ē $1356\frac{1}{4}$ plus leger que l'eau

p. 28. prop. 5\textsuperscript{e}

Il faut remarquer pour tout ce qui est icy dit, ou ailleurs, des rayons qui font le parallelisme, qu'il n'y a que ceux de la grandeur du soleil esgale a la grandeur du mi\add{roir} qui puissent garder vn parfait parallelisme iusques a l'infini, mais parce qu'il y en a \ill{} peu pō brusler, les autres qui viennent de chasque point du soleil, ne sont pas exacte\add{ment} paralleles, c'est pourquoy on ne peut faire ny miroirs qui bruslent, ny lunettes qui fac\add{ent}
\note{Au pied de la page, à droite, la réclame « voir les o… » ; la phrase se poursuit à la vue 38.}

\page{38}
\note{Page de droite (recto), chiffrée « 18 » à l'encre en haut à droite. Deux figures dans la marge de gauche, décrites en note.}

voir les obiets de mesme grandeur d'vne distance infinie. co\textsuperscript{e} l'on peut voir par ceste figure, dont, AHB, represente la surface du soleil parallele et esgale au verre CD, sur lequel si les seuls rayons paralleles AC, BD, et HG, auec tous les autres compris entre AB, et CD, venoient a tomber ilz se iroient tous aboutir au point G, supposé que le conuexe CM, fust hyperbolique, mais ilz ne pourroient brusler, sans l'ayde des autres rayons qui venans de chasque point du soleil, par exemple DA et B, s'estendent sur toute la surface du verre, et vont a peu pres se ioindre vers le point, G, lors qu'il n'est guieres esloigné du miroir, ou du verre, mais ceste infinité de rayons qui viennent de chasque point de l'obiect, sur chasque point du verre vont tousiours s'escartant de plus en plus de maniere qu'ilz se trouuent aussy esloignez les vns des autres co\textsuperscript{e} est le diametre du soleil, lors qu'ilz sont aussy esloignez du verre, co\textsuperscript{e} le verre du soleil, et partant ceux qui parlent de f\textsuperscript{e} des miroirs qui bruslent par ceste maniere iusques a l'infini, se trompent, toutesfois l'on peut f\textsuperscript{e} des lunettes de longue veüe suiuant la pre\add{miere} figure, ou mesme des miroirs qui brusleront aussy loing, ou qui feront voir \struck{voir} les obiets assez gros, ou clers, aussy loing, co\textsuperscript{e} iront les rayons non paralleles sans s'esloigner sensiblem\textsuperscript{t} du parallelisme. Mais il faut enfermer les deux glaces de miroir, ou des lunettes en vn tuyau de cuir, ou de charbon, ou d'autre matiere et le noircir en dedans, afin que la lumiere des costez n'entre point dedans, et que l'oeil posé en AD de la 1\textsuperscript{re} figure, ou AB de la seconde, ne reçoiue autre lumiere que celle des obietz OGF, ou MB, et pō ce subiet, le trou AD ne doibt pas estre plus grand que la paupiere de l'oeil, dont la largeur respondra a celle du diametre, ou de la largeur de la moindre parabole HB, et monstrera ce qui se perd des rayons. Or si la largeur de la grande parabole ē seulem\textsuperscript{t}. de 8 pouces, et celle de la petite d'vn demy pouce, les obietz paroistront plus clers, ou plus grands 256 fois qu'au seul oeil, moins les rayons qui se perdent dans l'espace d'vn demy pouce qui fait l'ouuerture AB.
\note{Suite de la remarque « p. 28. prop. 5\textsuperscript{e} » de la vue 37 (réclame « voir les o… »). Première figure, en haut de la marge : un rectangle H B en haut, D à droite, deux diagonales croisées ; dedans, un losange à sommet M, les points E et F sur l'axe, et G en bas, où convergent les lignes. Les lettres A et C ne sont pas lues sur la figure.}

voicy co\textsuperscript{e} il faut accommoder le tuyau

Ayant attaché la moindre parabole CD, a la grande FBEA aux points, F, E: ou a tels autres qu'on voudra, soit auec fil d'archal, lames de fer blanc, ou autrem\textsuperscript{t}, de sorte que les ataches facent perdre le moins de rayons que f\textsuperscript{e} se pourra, l'on enfermera le tout dans le tuyau GHIKL, lequel regardera les obiets, et si l'on veut on le fera en eslargissant en MN, afin qu'il descouure vne plus grande multitude d'obietz, ou plus grande partye du corps qu'on veut considerer, mais encore qu'il soit parallele sans s'eslargir il descouurira assez large, lors que les obietz seront fort esloignez, co\textsuperscript{e} il arriue aux estoiles, et aux planetes, et l'experience fera recognoistre plus de choses que nous n'en pouuons dire, par exemple combien il sera necessaire que le tuyau GIHK, soit long pō exclurre la lumiere laterale et cæt.
\note{Seconde figure, dans la marge : la grande parabole, sommet entre A et B en haut, passant par E et F ; la petite, C D, à l'intérieur ; le tuyau G H en haut, descendant vers L et K autour d'une ellipse (l'ouverture), et évasé jusqu'à N à droite.}

p. 30. lig. 20.

la petite partye du spherique qui doibt estre mise entre K et I, ou plustost en K, doibt auoir le mesme foyer que la p\textsuperscript{tie} de la grande sphere, co\textsuperscript{e} nous auons dit cy deuant des deux paraboles, dont le foyer doibt estre commun; et l'on a desia des miroirs spheriques fort bons, qui peuuent seruir a cest effect, en les perçant au milieu, et y adioustant vn autre petit miroir dans le conuexe mis vn peu plus haut que le milieu d'entre K et I, renuoyra les rayons receus sur le concaue \uncertain{Q O}, quasi paralleles, pō le tuyau qui doibt enfermer ces deux miroirs, il faut le f\textsuperscript{e} co\textsuperscript{e} j'ay dit pō les paraboles

p. 31. lig. 4.

Neantmoins l'obseruation faite de la plus grande hauteur ē plus certaine, co\textsuperscript{e} seroit celle du haut d'vne tour, ou d'vne montagne co\textsuperscript{e} il est dit cy dessous, a raison que le rayon qui frape l'eau de l'eleuation de 6 pieds la frape si fort de biays, qu'il est quasi parallele: mais en recompense la commodité ē plus grande, parce qu'il est aysé de mesurer bien iustem\textsuperscript{t} vne lieue a costé de l'estang sur lequel on fait l'obseruation, et l'obiect visible par exemple la chandelle, n'a pas besoing d'estre grosse, au lieu qu'és lieux fort esloignez, vne montagne qui sert d'obiect, ne paroist que co\textsuperscript{e} vn point.

\page{39}
\note{Page de gauche (verso). En tête, au-dessus de la figure, un « 81 » qui, retourné en miroir, se lit « 18 » : c'est le chiffre « 18 » du recto (vue 38) vu par transparence, non un nombre de cette page. Les fins de ligne passent dans le pli.}

p. 31. lig. 26.

Le grand diametre d'vne ellypse estant donné, l'on peut trouuer vn nombre qui explique toutes les lignes suyuantes de l'ellypse de sorte qu'elles seront toutes rationelles aud. grand diametre, a scauoir la ligne EF perpendicula\add{ire} a BA, la ligne AF, et le petit diametre CD, de sorte que la distance de \ill{} foyer au foyer H, soit plus grande que CD, et l'aire de l'ellypse plus grande que l'aire du cercle dont FH est le diametre, le nombre qui sert a cela pour 12 ellypses et non plus est 434788275897787529. Tel nombre d'ellypses qu'on voudra donner auec les conditions susd., on donnera vn nombre qui les exprimera, et tel nombre donné qu'on voudra, l'on peut dire a combien de tel\add{les} Ellypses il peut seruir, en faueur d'vn artisan qui demanderoit a f\textsuperscript{e} des ellypses ou tout cecy fut rationel pō f\textsuperscript{e} les lunettes dont parle la dioptrique de m\textsuperscript{r} des cartes, voyez la 9. page dans n\textsuperscript{re} exemplaire. l'analyse n'a peu seruir a ceste question iusques a present
\marginal{na}
\note{Figure dans la marge de gauche : une ellipse, grand axe AB horizontal, petit axe CD vertical ; le foyer F près de A et le foyer H de l'autre côté ; la perpendiculaire EF élevée en F jusqu'à l'ellipse ; une seconde perpendiculaire, repassée d'un trait gras, entre F et le centre. Le nombre à dix-huit chiffres est lu chiffre à chiffre sur une coupe agrandie ; le dernier chiffre, contre le pli, est le moins sûr. « na », dans la marge de gauche en regard de « des cartes, voyez la 9. page », est sans doute l'abréviation de « nota » ; la même marque se lit à la vue 36.}

p. 38. lig. 15.

voyez Galilée au 4\textsuperscript{e} dialogue, ou il tente de expliquer que ceste ligne ē parabolique, c'est en supposant que le mouuem\textsuperscript{t} parallele a l'horison ē tous\textsuperscript{rs} esgal en toutes ses parties ce qui repugne a l'experience: de sorte qu'il faut supposer qu'il n'y ayt point d'air, et que son mouuem\textsuperscript{t}. se face dans le vuyde. voyez depuis la 21 propo\textsuperscript{on} de la balistique iusques a la 32\textsuperscript{me}
\note{« de expliquer » : « que » est écrit au-dessus, entre « expliquer » et « ceste ».}

ligne 7 auant la fin

si le mouuem\textsuperscript{t}. horizontal ē tous\textsuperscript{rs} esgal en vistesse co\textsuperscript{e} veut Galilée, toute ceste suputation ē inutile, mais puis que la bale fait plus d'effect de pres que de loing, cela n'est pas vray, car si elle alloit tous\textsuperscript{rs}. esgalem\textsuperscript{t}. viste, elle feroit tous\textsuperscript{rs} mesme effect.

p. 39. lig. 19.

Galilée demonstre en son 4\textsuperscript{e} dialogue que la portée de 45 degrez ē 10000, lors que celle du p\textsuperscript{r} degré ē de 349, c'est a dire 28 fois moindre, mais si l'on y comprend \uncertain{l'empeschem\textsuperscript{t}} de l'air, il est fort esloigné de son conte; car la portée de 45 degrez dans nostre air, n'\add{est} guieres que triple que celle du p\textsuperscript{r} degré

p. 40 lig. 12 deuant la fin

Si le canon suit les portées d'arquebuse il est certain que tout ce qu'on dit des portées est faux, car l'experience fait voir que la portée de 45 degrez n'est pas quadruple de \add{celle} d'vn degré, ny de l'horizontale, lors que l'harquebuse n'est esleuée que de trois p\ill{} sur l'horizon. voyez le 3\textsuperscript{me} aduertissem\textsuperscript{t} qui suit

p. 42 au 3\textsuperscript{me} aduertissem\textsuperscript{t}. lig. 5.

\struck{aud} le dernier j\textsuperscript{r} de may \uncertain{1636}. sur la Foire a S\textsuperscript{t}. Clou.

La portée perpendiculaire en haut ē a l'horizontale de blanc en blanc co\textsuperscript{e} 14, a 5 \ill{} a celle de 4 degrez co\textsuperscript{e} 7 a 8 proxime.

La bale employe tous\textsuperscript{rs} vn temps a monter et descendre qui ē entre 22 et 26 second\add{es} lors qu'elle ē tirée par l'harquebuse

La grande portée sur Marne pres de Lagny ne nous a paru que de 330 toyses, ou enuir\add{on}
\note{La date se lit au début de la ligne, après un mot biffé, « aud » ; le premier chiffre de l'année est petit et peu formé, les trois autres sont nets.}

p. 43. lig. 7 auant la fin

les dernieres experiences monstrent que les flesches qui sont 8 secondes a monter et descendre, montent en 3 temps et descendent en 5 temps et partant qu'elles vont \ill{} haut que je n'ay dit.

p. 61. prop. 18.

Il y a 5 raisons a la marge de la 5\textsuperscript{e} obseruation pō prouuer que ceste opinion deduitte \ill{} elle est, est fausse, c'est a dire que sa preuue n'est pas bonne: J'adiouste icy, la 6\textsuperscript{me} a sca\add{voir} que faisant conceuoir la ligne FD, co\textsuperscript{e} vne balance dont le centre ē G, et mettant vn poi\add{ds au} point F, et vn autre au point E, qui pend du point D, il cherche le centre de grauité de \ill{} deux poids en la ligne EF, co\textsuperscript{e} s'ilz estoyent simplem\textsuperscript{t} ioints ensemble, ou il y a double erreur, car le poids E, pendant de D, en sorte que l'angle GDE, peut change\add{r} \ill{} mesme que la balance s'incline de part et d'autre; ne pese en ceste balance qu'autant q\ill{} tire le point D, et ainsy n'est opposé au point F, que suiuant la ligne FD, et non suiuant
\note{Figure dans la marge de gauche : un triangle F D A, le sommet A en bas ; G au milieu de FD, la perpendiculaire GA ; la droite FE coupant GA en H, E sur le côté DA. Au pied de la page, à droite, la réclame « la ligne » ; la phrase se poursuit à la vue 40.}

\page{40}
\note{Page de droite (recto), chiffrée « 19 » à l'encre en haut à droite.}

la ligne FE, Puis qu'il suposast que la ligne DE, fermement iointe a la ligne GD, en sorte que l'angle GDE ne peut se changer, toutesfois a cause du point G, qui estant le centre de la balance doit estre fixe, le centre de grauité des deux poids en F, et E, doibt estre tout autre que s'ilz estoyent considerez en vne balance.
\note{Fin de la remarque « p. 61. prop. 18. » de la vue 39 (réclame « la ligne »).}

Quant a ce qui ē de la verité de la question au fonds, voyez les lettres et discours de m\textsuperscript{r}. des Cartes la dessus
\marginal{na}
\note{« na » dans la marge de gauche, en regard de cet alinéa, comme aux vues 36 et 39.}

\marginal{\ill{} nouuelles \ill{}seruations \ill{}siques et mathematiques.}
\note{Titre courant dans la marge de gauche, sous un trait ; le début de chaque ligne est dans le pli. Sans doute « Des nouvelles observations physiques et mathematiques ». Les renvois qui suivent (p. 3, 4, 5) repartent d'une nouvelle pagination.}

p. 3. lig. 22.

voyez le 15 article du 1\textsuperscript{er} liure des pensées de Galilée, que j'ay mis en françois, afin de voir vne au\textsuperscript{re} maniere de peser l'air, qu'il trouue seulem\textsuperscript{t}. 400 fois plus leger que l'eau

lig. 12 deuant la fin

encores que la proportion de la vistesse des cheutes diminüe peu a peu, neantmoins l'on pretend qu'elles n'arriuent iamais a vn point d'egalité, duquel elles aprochent tous\textsuperscript{rs}, co\textsuperscript{e} les asymptotes des hyperboles, et de beaucoup d'autres figures, desquelles elles s'aprochent tous\textsuperscript{rs} sans y pouuoir paruenir

lig. 6. deuant la fin

voyez l'hydraulicopneumatique ou l'air ē pesé plus iustem\textsuperscript{t} et la marge du coroll. 8 de ceste proposition.

en la mesme ligne.

Ceste legereté d'air si grande ne semblera pas estrange, si l'on considere la quantité de la vapeur qui sort des Aeolipyles, comparée auec l'eau dud. Aeolipile

D'autres mesurent la densité de l'eau et de l'air, et des autres mylieux par le rayon rompu, ou reflechi, et disent que si la distance du point d'ou part le rayon d'auec la perpendiculaire ē seulem\textsuperscript{t}. vn tiers plus grande que la distance du rayon rompu d'auec la mesme perpendiculaire, l'eau sera seulem\textsuperscript{t}. d'vn tiers plus dense que l'air. le rayon tombant AC le rompu CD, la perpendiculaire FE, la distance dans l'air AB, dans l'eau ED, et pō fuir les dificultez, ilz nyent que la compression soit condensation, qu'ilz apellent vne qualité particuliere
\note{Figure dans la marge de gauche : un cercle de centre C, le diamètre vertical FE et le diamètre horizontal passant par C jusqu'à H ; B sur CF, la corde AB (A au bord gauche, non lettré sur la figure lue), le rayon incident de A en C et le rayon rompu CD, D en bas à droite, ED au pied.}

p. 4 au commencem\textsuperscript{t}.

J'ay trouué par la pompe Aeolipile que l'eau pese 1356 fois autant que l'air, ce qui reuient proxime a ceste proportion prise du mouuem\textsuperscript{t} de l'air et de l'eau comme de 1380 a 1.

lig. 15.

voyez la 33 propo\textsuperscript{on} des proportions de Cardan, ou il dit que les corps qui descendent d'esgale vistesse en des milieux differents, doiuent estre en raison doublée desd. milieux: d'ou il s'ensuit que supposé que le poids qui descend dans l'eau fust a celuy qui descend dans l'air, co\textsuperscript{e} 1600 a 1, tel qu'est presque le plomb a la vessie de poisson de mesme grosseur, l'eau peseroit 40 fois plus que l'air, puisque 40 ē la racine de 1600.
\note{« 1600 » à la fin de l'alinéa : écrit au-dessus d'un mot biffé.}

p. 5. au commencem\textsuperscript{t}

Tout ce qui est icy mis en doubte ē mainten\textsuperscript{t} resolu, car i'ay premierem\textsuperscript{t} trouué que les vistesses de la sortie de l'eau par vn mesme trou, sont en raison soubsdoublée des hauteurs des tuyaux, par exemple que les tuyaux ayant leurs hauteurs en raison de 4 a 1, la sortie de l'eau de celuy qui a 4 pieds de hauteur, ē a la sortie de l'eau du tuyau d'vn pied de hauteur, co\textsuperscript{e} 2 a 1, d'ou il s'ensuit qu'vn tuyau de 4 pieds abbaisse ou vuyde 2 fois autant son eau que celuy d'vn pied, et que celuy de 16 pieds de hauteur abbaisse ou vuide 4 fois plus viste son eau que celuy d'vn pied: mais il faut supposer qu'ilz soyent tous\textsuperscript{rs} tous pleins, et de mesme largeur.

2\textsuperscript{o}. J'ay trouué qu'vn tuyau de 4 pieds de hauteur, demeurant tous\textsuperscript{rs} plein d'eau, donne vne liure d'eau dans le temps de $23''$, et se desremplissant dans $25''$, d'ou il est aysé de conclurre combien le tuyau AB, de 12 ou d'autant de pieds qu'on voudra abbaisse son eau dans vn temps donné, et partant de quelle vistesse l'eau montera dans le canal CE; l'arge \struck{\ill{}} d'vne demye ligne, co\textsuperscript{e} ē le trou des tuyaux qui m'ont seruy d'essay.

3\textsuperscript{o}. le iet horizontal commenceant a 6 pouces de terre, ē de trois pieds de long, et celuy de 45 degrez ē de 6. pieds et entiers.

4\textsuperscript{o} son iet perpendiculaire ē de 3 pieds et $\frac{2}{3}$ pouces proxime, de sorte qu'il sort quasy vn quart de la hauteur du tuyau

lig. 9 auant la fin

Si l'eau monte d'autant plus viste perpendiculairem\textsuperscript{t} qu'elle sort plus viste, et que d'vn
\note{La phrase s'arrête au pied de la page, sans réclame ; elle se poursuit sans doute à la vue 41, hors de ce lot.}

\end{document}
